\section{Stelle di neutroni: Intro storica}

\begin{frame}{Brief History of NS}
    \begin{itemize}
        \item 1931 - Landau article on maximus mass of compact star: idea of compact star as giant nucleus
        \item 1932 - Chadwick discovers neutron
        \item 1934 - Article by Baade and Zwicky in which appeared for first time the Neutron stars: speculated aboud collapse of SN to neutron star
        \item 1934 - Oppenheimer, Volkoff studied neutron stars within General relativityi
        \item 1934 - Tolman ''Relativity, Thermodynamics and Cosmology'': fluid sphere in general relativity
        \item 1939 - TOV solution: OV numerically, T making assumption about metric. Maximum mass of star composed on gas of neutrons $0.7\msun{}$, using Tolman analytical solution shown than maxmass con be incresed not further of $1\msun$ including repulsion between neutrons.
        \item 1942 - Duyvendak: Chinise/Japonese records of ''guess star'' close to $\zeta$ Tauris in 1054 AD.
        \item 1942 - Mayall- Oort: Identified 1054 event with SN explosion and Crab nebula in Tauris constellation.
        \item 1958 - Harrison, Wheeler: Unified EOS connection dense matter from WD to NS. They used Empirical mass formula for nmuclei, the liquid drop model, and calculated EOS for mixture of nuclei, electrons and a free gas of neutron minimizing total energy in $\beta$-equilibrium. (Harrison et al 1965).
        \item 1967 - Cocke - Pulsar discovery. Association of star in middle of Crab nebula to Neutron star created by SN explosion.
        \item Cameron 1959 - Study Neutron Stars taking into account interactions between neutrons. Model of skyrme adapted for nuclear matter (Skyrme 1959). Maximum mass $2.0\msun$, radius \SI{8.1}{\kilo\meter}; predicted hyperons at higher densities.
        \item Salpeter 1960 - EOS including hyperons
        \item Ambartsumyan, Saakyan 1960 - Investigates appearence of hyperons in dense neutron star matter for Fermi gas: they found \PSigmaminus appears at density \SI{0.64}{\per\cubic\femto\meter} and \PLambda at \SI{1.27}{\per\cubic\femto\meter} which are 4 times and 8 times the saturation density $n_0=\SI{0.15}{\per\cubic\femto\meter}$
        \item Ambartsumyan, Saakyan 1962 - Mass-Radius relation of NS including hyperon for Fermi gas. EOS for nucleons and hyperons with interaction term proportional to number density squared - hyperons appear in core of massive NS; maximum mass in free case is $0.63\msun$ and radius \SI{11}{\kilo\meter}, in interacting case $1.0\msun$ and radius \SI{5.4}{\kilo\meter}.
\end{itemize}
\end{frame}

\section{Neutron Star Crust}

\subsection{Outer Crust}

\begin{frame}{Atmosphere}
    \begin{itemize}
        \item Transition atmosphere-solid crust: Fermi energy uguale a binding energy of last electron: $E_F=\frac{k_F^2}{2m_e}\approx E_b(Z)=\SI{13.6}{\ev}Z^2$
        \item Spessore atmosfera: da equilibrio idrostatico $\TDy{r}{P}\approx-\frac{P_c}{\Delta R}\approx-\frac{GM\rho_c}{R^2}$, quindi $\frac{\Delta R}{R}\approx \frac{R}{GM}\frac{2E_F}{5m_u}\frac{A}{Z}$; nel caso di stelle di neutroni lo spessore dell'atmosfera sar\'a $\frac{\Delta R}{R}\approx\SI{2.92e-8}(\frac{0.2}{C})ZA$: \SI{0.3}{\milli\meter} per atmosfera di H, \SI{10}{\cm} per atmosfera di Fe.
        \item Determine spectral properties: photon mean free path is smaller than atmospheric thickness.
        \item Atmospheric composition depends upon environment and formation conditions. For binary NS atmosphere will be H accreted from companion; for isolated NS moving in ISM composition will be mainly H; newly born NS will have composition related to material falling back from SN explosion dominated by heavy elements such Silicon.
    \end{itemize}
\end{frame}

\begin{frame}{BPS model} 
    \begin{itemize}
        \item Modello di Baym, Pethick, Sutherland (1971): reticolo di nuclei neutron-rich e gas di elettroni degeneri, assenza di neutroni liberi; $T=0$ equilibrio termodinamico/$\beta$-equilibrio; unica specie nucleare $(A,Z)$ per ogni strato; nuclei in reticolo coulombiano bcc; elettroni degeneri trattati come gas relativistico ideale; neutralit\'a elettrica $n_e=Zn_N$ dove $n_N$ \'e densit\'a dei nuclei; densit\'a barionica $n_b=An_N$; assenza neutroni liberi $n_n^{free}=0$.
        \item Per densit\'a superiori a \SI{1.42e4}{\gram\per\cubic\cm} l'energia di Fermi degli elettroni supera l'energia di legame dell'atomo pi\'u legato $^{56}Fe$ ($Z=26$). Sopra questa densit\'a gli elettroni formano un gas di Fermi libero e i nuclei minimizzano energia formando un lattice.
        \item contribuition of lattice to EOS  is small (see WD): pressure is dominated by degenerate elctron pressure. EOS of outer crust of NS is same of WD.
        \item Energy density is dominated by mass density of nuclei.
        \item Critical density for U R electrons: $E_F=k_F\approx m_e$ quindi $\rho_{rel}=m_un_N=\frac{m_uA}{3\pi^2Z}m_e^3\approx \SI{9.74e5}{\gram\per\cubic\cm}\frac{A}{Z}$.
        \item Change of composition: electron energy arriva a $m_n-m_p$ quindi si ha $\Pproton+Pelectron\to\Pneutron+\Pnue$. $E_F\approx m_n-m_p\approx\SI{1.29}{\mega\ev}$, $\rho_n=\frac{m_uA}{3\pi^2Z}(m_n-m_p)^3=\SI{1.57}{\gram\per\cubic\cm}\frac{A}{Z}$, con densit\'a numerica nucleoni $n_N=\frac{A}{Z}n_p=\frac{A}{Z}n_e$. For $^{56}Fe$ $\rho_n\approx\SI{3.37e7}{\gram\per\cubic\cm}$.
        \item Composizione determinata cercando il nucleo che minimizza energia tot. del sistema in funzione di $(A,Z)$
        \item Il potenziale termodinamico determinante \'e l'energia libera di Gibbs $G(P,T,N)$: \'e particolarmente semplice se consideriamo sistema a (P,T) costanti. Passiamo da U(S,V) con N conservato a G facendo trasformata di Legendre $S\to T$ e $V\to P$: $G=U+PV-TS$ e si ha $dG=VdP-SdT+\mu dN$; usando Eulero $U=-PV+TS+\mu N$ si ha $G=\mu N$: minimizzare G a fisso N coincide con minimizzare $\mu$.
        \item Per $T=0$ e usando le due espressioni ricavate al punto precedente per G si ha $\mu_b=\frac{\epsilon_{tot}+P_{tot}}{n_b}$ con $P_{tot}=n_b^2\PDy{n_b}{\frac{\epsilon_t}{n_b}}|_{(A,Z)}$ e visto che la densit\'a numerica di barioni \'e $n_b=An_i$, infine $P_t=P_e+\frac{1}{3}W_Ln_i$.
        \item EOS can be found minimizing G relto $(A,Z)$
    \end{itemize}
\end{frame}

\begin{frame}{Contribution of lattice to EOS is small: coulomb corrections}
    \begin{itemize}
        \item Spherical Wigner-Seitz cells with radius $r_0$ around nuclei.
        \item Charge neutrality: $Z=\frac{4\pi r_0^3}{3}n_e$.
        \item Coulomb self-energy of charged sphere: $E_{ee}=\int_0^{r_0}dq \frac{q(r)}{r}=\frac{3}{5}\frac{Z^2e^2}{r_0}$, using $q(r)=Ze \frac{r^3}{r_0^3}$
        \item Interaction electrons-ion: $E_i=\int_0^{r_0}dq \frac{Ze}{r}=-\frac{3}{2}\frac{Z^2e^2}{r_0}$ where charge of ion is $Ze$.
        \item $E_{coul}=E_{ee}+E_i=-\frac{9}{10}\frac{Z^2e^2}{r_0}$ hence the coulomb energy per electron from lattice is $\frac{E_{coul}}{Z}=-\frac{9}{10}Z\expy{\frac{2}{3}}e^2(\frac{4\pi}{3}n_e)\expy{\frac{1}{3}}$.
        \item Pressione dovuta al contributo coulombiano: $P_{coul}=n_e^2\TDy{n_e}{(E_{coul}/Z)}|_{(A,Z)}=-\frac{3}{10}Z\expy{\frac{2}{3}}e^2(\frac{4\pi}{3})\expy{\frac{1}{3}}n_e\expy{\frac{4}{3}}$. Politropa con $\Gamma=\frac{4}{3}$.
        \item Coulomb contribution is suppressed compared to Fermi pressure by 1-2 OOM: $Z\expy{\frac{2}{3}}e^2=\alpha Z\expy{\frac{2}{3}}=$ \num{0.01} for He - \num{0.03} for O. ($\alpha=e^2\approx \frac{1}{137}$).
    \end{itemize}
\end{frame}

\begin{frame}{Energia totale della cella di Wigner-Seitz: $\epsilon(A,Z,n_i)=\epsilon_e+n_i(W_L+W_i)$.}
    \begin{itemize}
        \item Considero reticolo bcc  (body centered cubic): ho 1 nucleo al centro della cella e un nucleo in ogni vertice condiviso fra altre 8 cella, quindi in media ho $1+8*\frac{1}{8}=2$ nuclei per ogni cella. Per semplicit\'a ricavo il raggio della sfera di wigner-Seitz attorno a ciascun nucleo: $\frac{2}{a^3}=\frac{1}{\frac{4\pi}{3}r_0}$ quindi $r_0\approx0.5a$.
        \item Neutralit\'a di carica $n_e=Zn_i$
        \item Total energy density of Wigner-Seitz cell: $\epsilon(A,Z,n_i)=\epsilon_e+n_i(W_L+W_i)$, dove $W_L$ is lattice energy and $W_i$ \'e l'energia degli ioni, ovvero la massa. $\epsilon_e$ \'e l'espressione standard per un gas di Fermi $\epsilon_e=\frac{g_e}{(2\pi)^3}\int d^3kE(k)=\frac{g_e}{2\pi^2}\int_0^{K_{F,e}}dkk^2\sqrt{k^2+m^2}$; il contributo del reticolo $W_L$ contiene energia coulombiana degli elettroni e degli ioni: $E_C=E_{ee}+E_i=-f\frac{Z^2e^2}{r_0}$, $f\approx \frac{9}{10}$.
        \item BCC has highest electrostatic potential so is the most stable configuration but difference between other structure is so small that we can have other lattice structure in outer crust of NS.
        \item La struttura del reticolo conta per le propriet\'a di trasporto: conducibilit\'a elettrica, conducibilit\'a termica e scattering elettronico dipendono da disposizione dei nuclei, fononi, vettori del reticolo reciproco, difetti.
        \item lattice poco importante per EO S: $\epsilon_t=\epsilon_i+\epsilon_e+\epsilon_L$, $\epsilon_L\ll\epsilon_i,\epsilon_e$ e anche $P_L=\frac{1}{3}\epsilon_L$ \'e piccola risp $P_e$ quindi i contributi del reticolo influiscono poco su $P(\epsilon)$.
        \item Il contributo de l reticolo \'e importante per determinare quali ioni compongono outer crust: a pressione fissata si minimizza $g(A,Z)=\frac{\epsilon+P}{n_b}$ e il contributo del reticolo \'e determinante per casi $g(A_1,Z_1)\approx g(A_2,Z_2)$.
    \end{itemize}
\end{frame}

\begin{frame}{Most important con tribution total energy of outer crust: nuclear mass $W_N$}
    \begin{itemize}
        \item $W_N=Zm_p+(A-Z)m_n+E_b(A,Z)$. Typical binding energy $-\frac{\SI{8}{\mega\ev}}{A}$; $^{56}Fe$ is the most stable: has lowest total energy per nucleon ($^{62}Ni$ has highest binding energy per nucleon).
        \item Missing nuclei from experimental data, which has to be considered in extreme conditions of NS,  has to be included via nuclear model calculations.
        \item Nuclei having magic number of protons/neutrons are particularly stable: numbers closed shell of protons/neutrons.
    \end{itemize}
\end{frame}

\begin{frame}{Nuclear shell model}
    \begin{figure}[!ht]\includegraphics[trim={0cm 0cm 0 0},angle=90,clip,keepaspectratio,width=0.95\textwidth]{Schaffner/magic-numbers-shell}\label{fig:Schaffner-magic-numbers-shell}\end{figure}
    Il pattern che emerge dal modello per composizione di Oter Crust of NS \'e analogo a quello dei numeri magici per chiusura shell protoni/neutroni del modello nucleare a shell.
    \begin{columns}[T]
        \begin{column}{0.5\textwidth}
            \begin{itemize}
                \item Posizioni shell chiuse
                \item Dimensione shell gaps
                \item Shell corrections alle masse
            \end{itemize}
        \end{column}
        \begin{column}{0.5\textwidth}
            \begin{itemize}
                \item $S_n$, $S_{2n}$ e drip-line
                \item Q-values delle catture $\beta$
                \item eventualmente livelli e rate di transizione
            \end{itemize}
        \end{column}
    \end{columns}
    
\end{frame}

\begin{frame}{''Modello BPS'' a P fissata in outer crust}
    \begin{itemize}
        \item Per ogni $(A,Z)$ la densit\'a di energia \'e $\epsilon=n_iW_i(A,Z)+\epsilon_e(n_e)+n_iW_L(A,Z,n_i)$ e pressione $P=P_e(n_e)+\frac{1}{3}n_iW_L$, $W_i$ masse nucleari sperimentali o ricavate da modelli nucleari. Fisso $P=P_1$ e cerco il minimo di g al variare di $(A,Z)$.
        \item Prendo nucleo candidato $(A_{11},Z_{11})$ con $n_i=\frac{n_e}{Z_{11}}$ quindi $P=P_e(n_e)+\frac{1}{3}n_iW_L(A_{11},Z_{11},n_e)$: risolvo $P_1=P(A_{11},Z_{11},n_e)$ rispetto a $n_e$ ottenendo $n_e(A_{11},Z_{11},P_1)$ e infine $n_i=\frac{n_e}{Z}$, $n_b=\frac{A}{Z}n_e$.
        \item Calcolo potenziale di Gibbs per barione. A $T=0$ ho $g=\frac{G}{N_b}=\mu_b=\frac{\epsilon+P}{n_b}$ infatti sostituendo $U=TS-PV+\sum_a\mu_aN_a$ in $G=U-TS+PV$ quindi $\sum_a\mu_aN_a=\mu_BN_B+\mu_QN_Q+\mu_LN_L$, inoltre la neutralit\'a impone $N_Q=0$ e per $\nu s$ non intrappolati $\mu_L=0$, infine $g=\frac{G}{n_B}=\frac{V(\epsilon+P)}{n_BV}$ e riscriviamo $g(A,Z;P)=\frac{1}{A}[W_i(A,Z)+Z\mu_e+\frac{4}{3}W_L]$
        \item Determinare minimo di $g$ al variare di $(A,Z)$
        \item Aumentare la pressione: $P_2>P_1$.
        \item Per ogni valore della pressione ho determinato $n_b=\frac{A}{Z}n_e$ e $\rho=m_un_b$.
        \item \xaumenta{P}, \xdiminuisce{\frac{Z}{A}}, compare stato con neutroni liberi energeticamente favorito: esco da modello BPS.
    \end{itemize}
\end{frame}

\begin{frame}{Sequence of nuclei jump discontinuosly: first-order phase transition.}
    \begin{itemize}
        \item Se la sequenza di nuclei passa ad uno pi\'u stabile ho discontinuit\'a nella densit\'a di energia dovuta alla differenza nelle binding energy - first order phase transition. A $T=0$ $\epsilon=n_b\mu_b-P$ e $P, \mu_b$ sono uguali nelle due fasi ma $n_b$ cambia: $\Delta\epsilon=\mu_t\Delta n_b\neq0$.
        \item Costruzione di Maxwell - Al punto di transizione ho $P_{\alpha}=P_{\beta}=P_t$, $\mu_{\alpha}=\mu_{\beta}$: abbiamo transizione da composizione $\alpha$ a $\beta$ a $P$ costante e dato che $P(r)$ nelle stella \'e strettamente monotono si ha una superficie di discontinuit\'a per $(A,Z)$, $n_b$, $\epsilon$.
    \end{itemize}
\end{frame}

\begin{frame}{Modelli nucleari}
    \begin{itemize}
        \item Parametrized mass model: historical originated from liquid drop model (Weizsacker 1935, Bethe and Bacher 1936). Bethe-Weizsacker mass formula
            \begin{equation*}
                B(A,Z)=a_vA+a_{surf}A\expy{\frac{2}{3}}+a_CZ^2A\expy{-\frac{1}{3}}+a_{sym}\frac{(N-Z)^2}{A}+\delta(A)
            \end{equation*}
            where $N=A-Z$, nuclear radius $R=r_0A\expy{\frac{1}{3}}$, $r_0\approx\SIrange{1.0}{1.2}{\femto\meter}$: first term ($\propto R^3$) is the volume term, second term ($\propto R^2$) is the surface term, third term ($\propto \frac{Z^2}{R}$) is proportional to coulomb potential of charges sphere with charge Z e radius R, the fourth term is symmetry term describing isospin dependence of nuclear forces, the last term stands for effect of odd number of protons/neutrons: $\delta(A)=\pm a_{pair}A\expy{-\frac{1}{2}}$ resp for EE and OO numbers of pn and $\delta(A)=0$ if A is odd. Fits the measured nuclear masses (see chowdhury, basu 2006) give coeff: $a_v=-\SI{16}{\mega\ev}$, $a_s=\SI{18}{\mega\ev}$, $a_C=\SI{0.7}{\mega\ev}$, $a_{sym}=\SI{23}{\mega\ev}$, $a_{pair}=\SI{-12}{\mega\ev}$.
        \item Skyrme models: Parametric ansatz for nucleon mean effective potential within \schr{}: potential terms have zero range and can be given in terms of proton/neutron density and gradient thereof. See ''The Skyrme Interaction in finite nuclei and nuclear matter'', ''Skyrme Crystals, Nuclear Matter and Compact Stars'', ''The Skyrme Interaction in finite nuclei and nuclear matter'', ''A unified equation of state of dense matter and neutron star structure'', ''Physics of Neutron Star Crusts''.
        \item Relativistic Mean Field Models - Motivated from meson exchange picture of nuclear forces as proposed by Yukawa and is originated from work of Walecka 1974. Il potenziale del nucleone, nell'equazione di Dirac, \'e mediato dal compo mesonico, a loro volta soluzioni dell'equazione di Klein-Gordon, con sorgente campo mesonico costituito da densit\'a dei nucleoni.
    \end{itemize}
\end{frame}

\begin{frame}{Sequence of nuclei in Outer Crust of NS}
    \begin{columns}[T]
        \begin{column}{0.55\textwidth}
    
    \begin{figure}[!ht]\includegraphics[trim={0cm 0cm 0 0},clip,keepaspectratio,width=0.95\textwidth]{Schaffner/sequence-nuclei-crust-NS}\label{fig:Schaffner/sequence-nuclei-crust-NS}\end{figure}
            
        \end{column}
        \begin{column}{0.45\textwidth}
    \begin{itemize}
        \item Sequence of nuclei in crust of NS: sequence stars at most stable nucleus $^{56}Fe$.
        \item All'aumentare della densit\'a e quindi della densit\'a barionica aumenta anche la densit\'a elettronica per la neutralit\'a di carica (vedi aumento potenziale chimico nelle prime due colonne).
        \item Si raggiunge densit\'a per cui conviene convertire protone e elettrone in neutrone: avremo nuclei neutron rich. Per densit\'a $\rho\approx\SI{8e6}{\gram\per\cubic\cm}$ $^{62}Ni$ \'e energeticamente favorito.
    \end{itemize}
            
        \end{column}
    \end{columns}
            \begin{itemize}
                \item $^{62}Ni$ has magic number $Z=28$; at densities of $\rho\approx\SI{3e9}{\gram\per\cubic\cm}$ appaiono nuclei con numero magico $N=50$.
                \item $^{80}Zn$ is the last nucleus with measured BE. Beyond densities $\rho\ge\SI{6e10}{\gram\per\cubic\cm}$ nuclei are so neutron rich that their masses aren't measures.
                \item Relativistic mean field parameter set TMA predict doubly magic $^{78}Ni$ and that sequence continue along neutron magic number $N=82$
            \end{itemize}
    \begin{itemize}
        \item $^{116}Se$ is predicted to be last nucleus in OC at $\rho\approx\SI{4.55e11}{\gram\per\cubic\cm}$. Here baryon chemical potential have reached neutron mass $\mu_b=m_n=\SI{939.57}{\mega\ev}$: free neutron can be added to the system not bounded to nuclei. ($\mu_b=(\PDy{N_b}{G})_{P,T}$ rappresenta il costo in Gibbs per aggiungere barione al sistema: finch\'e $\mu_b<m_n$ un neutrone libero \'e pi\'u costoso di un barione legato; a densit\'a di neutroni liberi infinitesima un neutrone libero ha $p_{F,n}\to0$ quindi $\mu_n^{free}\to m_nc^2$ e quindi la presenza di un neutrone libero non aumenta Gibbs del sistema, ie $dG_{nuclei}=dG_{free}$).
        \item Neutron drip line marks begins of inner Crust.
    \end{itemize}
\end{frame}

\begin{frame}{Sequence of nuclei in nuclear chart}
    \begin{figure}[!ht]
        \includegraphics[trim={0cm 0cm 0 0},clip,keepaspectratio,width=0.5\textwidth]{Schaffner/nuclearchart-sequence-nuclei-crust-NS}\label{fig:Schaffner-nuclearchart-sequence-nuclei-crust-NS}
    \end{figure}
    \begin{itemize}
        \item Tipicamente in OC abbiamo $Z\approx\numrange{26}{46}$
        \item La carica dei nuclei \'e determinata dal beta-equilibrio: per densit\'a superiori a \SI{6e11}{\gram\per\cubic\cm} si hanno piccoli discostamenti di carica nucleare dovuti alle differenze nei modelli nucleari.
    \end{itemize}
\end{frame}

\begin{frame}{EOS of NS Outer Crust.}
            %lbrt
            \begin{figure}[!ht]\includegraphics[trim={0.4cm 0.4cm 0.3cm 0.5cm},clip,keepaspectratio,width=0.6\textwidth]{Schaffner/EOS-NS-OC}\label{fig:Schaffner-EOS-NS-OC}\end{figure}
    \begin{itemize}
        \item EOS don't shows significative difference for different nuclear models used: nuclear binding energy per nucleone circa \SI{8}{\mega\ev}, massa nucleoni approx \SI{1}{\giga\ev}.
    \end{itemize}
\end{frame}

\begin{frame}{EOS of NS Outer Crust: thickness and mass}
    \begin{align*}
        &\TDy{r}{P(r)}=-\frac{Gm_r(r)\epsilon(r)}{r^2}(1+\frac{P(r)}{\epsilon(r)})(1+\frac{4\pi r^3P(r)}{m_r(r)})\invers{(1-\frac{2Gm_r(r)}{r})}\tag{TOV:HE+GR terms}
    \end{align*}
            %lbrt
    \begin{columns}[T]
        \begin{column}{0.5\textwidth}
            \begin{figure}[!ht]\includegraphics[trim={0.5cm 0.5cm 0.2cm 0.5cm},clip,keepaspectratio,width=0.95\textwidth]{Schaffner/mass-thickness-OC}\label{fig:Schaffner-mass-thickness-OC}\end{figure}
        \end{column}
        \begin{column}{0.5\textwidth}
    \begin{align*}
        &\Rightarrow -\frac{\Delta P}{\Delta R}\approx \frac{P_b}{\Delta R}\approx \frac{GM\epsilon_b}{R^2(1-\frac{2GM}{R})}=\frac{C}{1-2C}\frac{\epsilon_b}{R}    \end{align*}
    \begin{itemize}
        \item As EOS of OC is well determined: solve TOV with modified initial conditions. We assume $\Delta R\ll R$, mass of OC negligible relative to $M_{NS}$, $P_b\ll\epsilon_b$: $\frac{\Delta R}{R}\approx \frac{1-2C}{C}\frac{P_b}{\epsilon_b}$, con compactness $C=\frac{GM}{R}\approx\numrange{0.2}{0.3}$, riscriviamo $\frac{\Delta R}{R}\approx \frac{1-2C}{C}\frac{P_b}{\epsilon_b}$ (GR correction squeeze the crust rel. to newtonian).
    \end{itemize}
        \end{column}
    \end{columns}
    \begin{itemize}
        \item Pressure to energy density ratio. UR \Pelectron: $\mu_e\approx k_{F,e}=(3\pi^2n_e)\expy{\frac{1}{3}}$
            \begin{align*}
                &\frac{P_{e,rel}}{\rho_i}=\frac{(3\pi^2)\expy{\frac{1}{3}}}{4}\frac{n_e\expy{\frac{4}{3}}}{m_uAn_i}=\frac{(3\pi^2)\expy{\frac{1}{3}}}{4}\frac{Zn_e\expy{\frac{1}{3}}}{Am_u}=\frac{Z\mu_e}{Am_u}\Rightarrow \frac{P_b}{\epsilon_b}\approx\num{3e-3}\tag{i.b.: $\frac{Z}{A}=\frac{34}{82}\approx0.3$, $\mu_e\approx\SI{26.29}{\mega\ev}$}\\
                &\Delta R_{OC}=\frac{1-2C}{C}\frac{Z\mu_e}{4Am_u} R\Rightarrow \approx \SI{40}{\meter} (2\msun{}, \SI{10}{\kilo\meter}, C=0.3)\ \SI{200}{\meter} (1.4\msun{}, \SI{15}{\kilo\meter}, C=0.15)\\
                &\Delta M=4\pi\int_0^{\epsilon_b}drr^2\epsilon(r)\approx 4\pi R^2\rho_b\Delta R\approx\num{2.5e-4}\msun{}(\frac{R}{\SI{10}{\kilo\meter}})^2(\frac{\rho_b}{\SI{4e11}{\gram\per\cubic\cm}})(\frac{\Delta R}{\SI{1}{\kilo\meter}})
            \end{align*}
    \end{itemize}
\end{frame}

\subsection{Inner Crust}

\begin{frame}{Inner Crust Phenomenology}
    \begin{itemize}
        \item Inner crust begins at neutron-drip-line, when neutron start being ejected out of nuclei, forming a separate liquid. Dripline is where the baryon chemical potential equals vacuum neutron mass: neutron drip density is model dependent $\rho_{drip}\approx\SIrange{4e11}{5e11}{\gram\per\cubic\cm}$
        \item Relation between baryon density and baryon chemical potential is not unique; baryon chemical potential is energy needed to add/subtract baryon (neutron) from the medium, which depends on potential energy the baryon feels in the medium. For outer crust potential energy of baryon is just binding energy of the nucleus, so the binding energy of last nucleus in the sequence has to be known with precision to determine neutron drip-density.
    \end{itemize}
            \begin{columns}[T]
                \begin{column}{0.45\textwidth}
                    \begin{align*}
                        &g(A,Z;P)=\frac{\epsilon+P}{n_b}\tag{$T=0$}\\
                        &g(A,Z;P)=\frac{1}{A}[W_N(A,Z)+Z\mu_e+\frac{4}{3}W_L]\\
                        &M(A,Z)=Zm_p+(A-Z)m_n-B(A,Z)
                    \end{align*}
                \end{column}
                \begin{column}{0.55\textwidth}
                    \begin{align*}
                        &S_n(A,Z)=M(A-1,Z)c^2+m_nc^2-M(A,Z)c^2\tag{vacuum}\\
                        &S_n<0\tag{emissione neutrone}\\
                        &g(A,Z;P)=\mu_b\tag{OC}
                    \end{align*}
                \end{column}
            \end{columns}
            \begin{itemize}
                \item La struttura della inner crust \'e simile a OC nel fatto che protoni sono all'interno degli ioni organizzati in reticolo circondati da gas di elettroni per bilanciare la carica ma i neutroni sono presenti anche fuori dai nuclei formando un liquido separato.
                \item Il fluido di neutroni \'e superfluido - T. di Cooper (Basis of BCS): degenerate fermions will form bound pairs if there is an attractive potential between them - just one attractive channel of neutron-neutron interaction will make neutron matter superfluid even if the sum of all channels is repulsive. Nucleon scattering experiment: partial wave $^1S_0$ is attractive ($^{2s+1}L_J$), extrapolation of nuclear interaction model to neutron fluid predicts \SI{1}{\mega\ev} attractive potential (at higher density other channels can turn attractive).
                \item Neutron superfluid is corotating (ms-pulsar: rapidly rotating). Rotating superfluid form vortices: quantized angular number momentum trasport.
                    Rotating superfluid is irrotational everywhere except for vortices core. $m\oint v_s\cdot dl=sh$; for uniform rotation $\vec{v}=\vecp{\Omega}{r}$ quindi $\nabla\wedge\vec{v}=\nabla\wedge\vecp{\Omega}{r}=2\vec{\Omega}$, vorticit\'a di un fluido classico in rotazione uniforme.
            \end{itemize}
\end{frame}

\begin{frame}{Vortex structure in inner crust}
    \begin{itemize}
        \item Vorticit\'a di fluido classico in rotazione uniforme $\nabla\wedge\vec{v}=2\vec{\Omega}$; lontano dai vortici $\nabla\wedge\vec{v}_s=0$.
        \item Supponiamo superficie A ortogonale a $\vec{\Omega}$ contenga molti vortici $N_V$ con $s=1$: la circolazione di A \'e la somma di tutti i vortici racchiusi $\oint_{\partial A}\vec{v}_s\cdot d\vec[\mathbf]l=N_V \frac{h}{m}$.
        \item T. Stokes $\oint_{\partial A}\vec{v}_s\cdot d\vec{l}=\int_A\nabla\wedge\vec{v}_s\,dS$: $m\oint_{\partial A}\vec{v}_s\cdot d\vec{l}=N_V \frac{h}{m}=2\Omega A$ quindi $\frac{N_V}{A}=\frac{2\Omega m}{h}$ (Relazione di Feynman per densit\'a dei vortici). Densit\'a momento angolare rispetto asse del vortice $\vec[\mathbf]{l}(\vec{r})=\vec{r}\wedge\rho_s\vec{v}_s$ quindi per vortice assisimmetrico $v_{\phi}(\vec{r})=\frac{s\kappa}{2\pi r}=\frac{s\hbar}{mr}$ dove s \'e il winding number, e $l_z(r)=\rho_srv_{\phi}=\frac{\rho_ss\hbar}{m}$.
        \item I vortici formano reticolo bidimensionale ortogonale all'asse di rotazione: in termini del periodo P della pulsar $\frac{N_V}{A}=\frac{8\pi m_n}{hP}=\SI{6.35e7}{\per\square\meter}\frac{\si{\second}}{P}$. The space between vortices is much smaller than NS size and much bigger than microscopic lattice scale of outer crust.
        \item Pulsar Glitches: Pulsar suddenly rotates faster. Angular momentum transport by vortices. Schematicamente: Spin-down elettromagnetico della crosta, vortici pinned (ostacolati nel moto dalla presenza del reticolo di nuclei), $\Omega_s-\Omega_c$ cresce finch\'e non si ha unpin di molti vortici che si muovono verso l'esterno.
    \end{itemize}
\end{frame}

\begin{frame}{young NS cooling: neutrino emission, meccanismo PBF per $T<T_c$}
    \begin{itemize}
        \item Young NS cool emitting neutrinos. La superfluidit\'a influisce sul raffreddamento di NS attraverso emissione di neutrini. Neutron pairs of bound neutrons have to be broken: per $T<T_c$ in inner crust prevales pairing $^1S_0$ between neutron which produce coherent mixing of many pairs $(k,-k)$ producing BCS condensate. The probability for breaking neutron pair at low $T\neq0$ is suppressed by Boltzmann factor $\exp{-\frac{2\Delta}{kT}}$, which also suppresses neutrino emission rate. Nel superfluido neutronico le eccitazioni elementari sono quasi-particelle di Bogoliubov con $E_k=\sqrt{\xi_k^2+\Delta^2}$, rompere una coppia significa creare due quasi-particelle: condensato $\to\gkupd{k}+\gmkdn{k}$ con gap $2\Delta$. A T finito esiste equilibrio BCS-eccitazioni: $\ketBCS\underset{\text{pair formation}}{\overset{\text{pair breaking}}{\rightleftarrows}}\gkupd{k}\gmkdnd{k}\ketBCS$, $0<T<T_c$. \item Osservazioni della curva di raffreddamento forniscono informazioni su $T_c$ tramite $\Delta(T=0)\approx\num{1.764}kT_c$.
        \item Il \(^{1}S_0\) neutronico della crosta interna \'e  molto importante per la superfluidit\'a e per la dinamica dei vortici, ma nelle curve globali di cooling il gap neutronico del core, soprattutto \(^{3}P_2\!-\!{}^{3}F_2\), \'e generalmente quello che ha generato i vincoli osservativi pi\'u discussi attraverso Cas A e il PBF.
        \item Il pairing ha due effetti fondamentali: sopprime \(C_V\) e molti processi neutrino per \(T\ll T_c\), ma introduce l'emissione PBF quando \(T\) attraversa \(T_c\).
        \item Nel 2011 Page e indipendentemente Shternin interpretarono il rapido calo della temperatura superficiale della giovane NS di Cas A come l'accensione del PBF quando il core neutronico attraversa \(T_c\) per il pairing \(^{3}P_2\). Page et al ottennero una scala \(T_c\sim5\times10^8\,{\rm K}\); Shternin et al. richiedevano un profilo \(T_{cn}(\rho)\) ampio con massimo circa \(7\!-\!9\times10^8\,{\rm K}\). Entrambi richiedevano inoltre protoni gi\'a superconducting, cos\'i da sopprimere parte del cooling precedente.
    \end{itemize}
\end{frame}

\begin{frame}{Cooling delle neutron star e pairing nucleonico}
\small
\begin{itemize}
  \item \textbf{Yakovlev \& Pethick (2004), \emph{Neutron Star Cooling}.}
  Review classica sul raffreddamento delle neutron star:
  processi neutrino, superfluidità, calore specifico e confronto
  fra curve di cooling e osservazioni.
  \item \textbf{Potekhin, Pons \& Page (2015),
  \emph{Neutron Stars---Cooling and Transport}.}
  Review moderna dell'evoluzione termica, trasporto di calore,
  struttura di core/crosta/envelope e interpretazione
  della radiazione termica osservata.
  \item \textbf{Sedrakian \& Clark (2019),
  \emph{Superfluidity in Nuclear Systems and Neutron Stars}.}
  Review microscopica del pairing nucleonico:
  neutroni $^1S_0$ nella crosta, neutroni $^3P_2$--$^3F_2$
  nel core, gap superfluidi, interazioni deboli e vortici.
  \item \textbf{Page et al. (2011),
  \emph{Rapid Cooling of the Neutron Star in Cassiopeia A
  Triggered by Neutron Superfluidity in Dense Matter}.}
  Interpretazione del rapido cooling di Cas A come onset
  del pairing neutronico tripletto e conseguente emissione PBF.
  \item \textbf{Shternin et al. (2011),
  \emph{Cooling neutron star in the Cassiopeia A supernova remnant:
  evidence for superfluidity in the core}.}
  Interpretazione indipendente di Cas A; vincoli sul profilo
  $T_{c,n}(\rho)$ e sul massimo della temperatura critica
  del pairing neutronico nel core.
  \item \textbf{Beloin et al. (2018),
  \emph{Constraining Superfluidity in Dense Matter from the Cooling
  of Isolated Neutron Stars}.}
  Inferenza statistica di $T_c$ neutronico e protonico
  da un insieme di neutron star isolate nel quadro del
  \emph{minimal cooling}.
\end{itemize}
\end{frame}

\begin{frame}{Superfluidit\'a: BCS 101}
    A $T=0$ un gas di fermioni non interagenti riempe tutti gli stati fino $E_F$, $\xi_k=\epsilon_k-\mu$ ($\xi_k<0$ per $\epsilon_k<E_F$): il pairing \'e dominato da stati con $|\xi_k|\ll E_F$. Attrazione debole vicino $E_F$ produce stato legato: Cooper considera due fermioni sopra Fermi's sea inerte - la coppia pi\'u semplice \'e formata da stati time-reversed ($\vec[\mathbf]{k}\uparrow,\vec[\mathbf]{k}\downarrow$), momento totale nullo e pairing in S-wave. Il fermi sea non \'e il GS in presenza di attrazione. (In BCS si considerano si stati sopra, sia stati sotto $E_F$, in cooper solo stati sopra). BCS estende da singola coppia a intero sistema: $\ket{BCS}=\prod_{\mathbf{k}}\left(u_k + v_k \pairkd\right)\ket{0}$, per ogni pair-state k, $u_k^2$ \'e la probabilit\'a che la coppia sia assente, $v_k^2$ la probabilit\'a che sia presente, con $u_k^2+v_k^2=1$. Il parametro di pairing $\kappa_k=\avg{\cmkd}{\cku}$, nel ground-state BCS $\kappa_k=u_kv_k$, \'e l'ampiezza di pairing dello stato di coppia l e attraverso $V_{kl}$ contribuisce al campo efficace sentito da coppia k: il campo di pairing su k $\Delta_k=-\sum_lV_{kl}\kappa_l$. Possiamo ottenere l'espressione precedente considerando $H_{\mathrm{int}}^{{}^1S_0}=\sum_{kl}\underbrace{\left\langle k;{}^1S_0\middle|V_{nn}\middle|l;{}^1S_0\right\rangle}_{V_{kl}}\,\underbrace{b_k^\dagger b_l}_{\text{pair scattering}}$, dove $\bpair{l}=\cmkdn{l}\ckup{l}$. Introduco mean-field approx $\bpair{k}\approx\overbrace{\avg{\bpair{k}}}^{\kappa_k}+\delta\bpair{k}$, dinamica collettiva complicata→ogni grado di libert\'a vede il campo medio degli altri, con $\delta\bpaird{k}=\bpaird{k}-\avg*{\bpaird{k}}$, $\delta\bpair{l}=\bpair{l}-\avg*{\bpair{l}}$ quindi $\bpaird{k}\bpair{l}=\avg*{\bpaird{k}}\avg*{\bpair{l}}+\delta\bpaird{k}\,\avg*{\bpair{l}}+\avg*{\bpaird{k}}\,\delta\bpair{l}+\delta\bpaird{k}\,\delta\bpair{l}$ quindi $H_{int}^{MF}=-\sum_k(\Delta_k\bpaird{k}+\Delta_k^*\bpair{k})+\sum_k\kappa_k^*\Delta_k$, a $H_{int}^{MF}$ aggiungiamo $\sum_{k\sigma}\xi_k\,c^\dagger_{k\sigma}c_{k\sigma}=H_0-\mu\hat N$ con $\xi_k=\epsilon_k-\mu$ e nel modello di particelle libere $\epsilon_k=\frac{\hbar^2k^2}{2m}$, termine a una particella, per avere $H^{MF}$: il primo termine risulta da $\sum_{kl}V_{kl}\bpaird{k}\kappa_l=\sum_k\bpaird{k}(\sum_lV_{kl}\kappa_l)=-\sum_k\Delta_k\bpaird{k}$; per il secondo termine considero interazione hermitiana $V_{kl}=V_{lk}^*$, quindi $\Delta_k^*=-\sum_lV_{lk}^*\kappa_k^*=-\sum_lV_{kl}\kappa_k^*$, nel modello BCS semplice $\Delta_k=\Delta$. Nel normale gas di Fermi l'energia rispetto a $\mu$ \'e $|\xi_k|$ mentre ne lsuperfluido BCS $E_k=\sqrt{\Delta^2+\xi_k^2}$; alla superficie di Fermi $\xi_k=0$ quindi $E_{min}=\Delta$. Le eccitazioni sono combinazioni particle/hole $\gamma_k=u_kc_k-v_kc_{-k}^{\dag}$, quasi-particelle di Bogoliubov: si ha che $u_k^2=\frac{1}{2}(1+\frac{\xi_k}{E_k})$, $v_k^2=\frac{1}{2}(1-\frac{\xi_k}{E_k})$ con $u_kv_k=\frac{\Delta}{2E_k}$, vicino $k_F$ $u_k^2\approx v_k^2\approx \frac{1}{2}$, quindi il mixing particle-hole \'e massimo. Il pairing smussa l'occupazione attorno a $E_F$, quindi tende ad aumentare l'energia cinetica, ma abbassa maggiormente l'energia di interazione. Il risultato netto \'e $U_s-U_n = -\frac{1}{2}N(0)\Delta^2$ nel weak-coupling BCS semplice.
\end{frame}

\begin{frame}{Inner Crust EOS}
    \begin{itemize}
        \item Reticolo di celle di Wigner–Seitz, ciascuna contenente un cluster nucleare ricco di neutroni immerso in un gas di neutroni ''drippati'' e in un fondo elettronico: Negele e Vautherin (1973) trattano ogni cella come una sfera di raggio \(R_{\rm WS}\) e risolvono in modo self-consistent Hartree–Fock il problema nucleare nella cella: cercano le densit\'a $n_p(r)$, $n_n(r)$ che minimizzano l'energia per data densit\'a media della materia, variando anche composizione $(N,Z)$ e dimensione della cella.
        \item I protoni hanno picco di densit\'a al centro della cella, i neutroni hanno picco al centro della cella e profilo piatto non nullo verso l'esterno; all'aumentare della densit\'a aumenta la densit\'a neutronica nella parte di piatta, infine si restringe anche la cella e il profilo neutronico della densit\'a \'e costante con piccolo incremento centrale.
        \item Strong shell effects for protons: for lower densities number of protons inside the cell is $Z=40$, at higher densities switches to $Z=50$: a magic number for the shell model without a strong spin-orbit splitting and to a closed subshell for the case of strong spin-orbit splitting; proton numbers staying around \numrange{40}{50} is a balance between the surface energy and the Coulomb energy of the Wigner-Seitz cell; the total neutron number inside a cell reaches $N=1800$.
        \item Inner crust stops at $0.5n_0\approx\SI{0.075}{\per\cubic\femto\meter}$ ($0.5\epsilon_0=0.5m_nn_0\approx\SI{70}{\mega\ev\per\cubic\femto\meter}$), lattice size shrinks from \SI{100}{\femto\meter} to \SI{30}{\femto\meter}.
        \item For higher densities, denpending upon nuclear model, we reach outer core consisting of liquid neutrons with small admixture of protons, electrons; lattice disappears.
    \end{itemize}
\end{frame}

\begin{frame}{Hartree--Fock autoconsistente nella cella di Wigner--Seitz}

\small
\begin{itemize}
  \item Si approssima una cella della crosta interna con una sfera di raggio
  \(R_{\mathrm{WS}}\), assumendo simmetria sferica:
  \[
    n_n(\mathbf r)\to n_n(r),
    \qquad
    n_p(\mathbf r)\to n_p(r).
  \]

  \item Nel metodo Hartree--Fock il ground state è approssimato da un determinante
  di Slater e gli orbitali soddisfano
  \[
    h_{\mathrm{HF}}[n_n,n_p]\phi_i=\epsilon_i\phi_i .
  \]

  \item L'autoconsistenza consiste nel richiedere che il campo medio costruito
  dalle densità produca, risolvendo le equazioni HF, le stesse densità:
  \[
    n_q(r)=\sum_{i\in q}^{\mathrm{occ}}|\phi_i(r)|^2,
    \qquad q=n,p .
  \]

  \item Schema iterativo:
  \[
  n_q^{(0)}(r)
  \rightarrow
  h_{\mathrm{HF}}[n_q^{(0)}]
  \rightarrow
  \phi_i
  \rightarrow
  n_q^{(1)}(r)
  \rightarrow
  \cdots
  \rightarrow
  n_q^{(m+1)}(r)\simeq n_q^{(m)}(r).
  \]

  \item La configurazione della cella si ottiene minimizzando l'energia totale
  rispetto a
  \[
    N,\qquad Z,\qquad R_{\mathrm{WS}} .
  \]
\end{itemize}

\end{frame}

\begin{frame}{Stima spessore crosta}
    \begin{itemize}
        \item For inner pressure we take pressure of NR neutrons, instead of pressure of UR electron as in outer crust. At inner boundary of inner crust (gas degenere NR, $T=0$)
            \begin{columns}[T]
                \begin{column}{0.42\textwidth}
            \begin{align*}
                &\frac{P_b}{\epsilon_b}|_{crust}\approx\xi \frac{P_n^{NR}}{\rho_n}=\frac{\xi}{5}(3\pi^2)\expy{\frac{2}{3}}\frac{n_n\expy{\frac{2}{3}}}{m_n^2}=\frac{\xi}{5}\frac{k_{F,n}^2}{m_n^2}\\
                &E_{F,n}=\frac{k_{F,n}^2}{2m_n}\Rightarrow \frac{P_b}{\epsilon_b}\approx \frac{2\xi}{5}\frac{E_{F,n}}{m_n}\\
                &\approx \frac{0.4}{5}(\frac{257}{939})^2\approx\num{6e-3}
            \end{align*}
                \end{column}
                \begin{column}{0.58\textwidth}
                    $P_b=\xi P_n^{NR}$, $\xi=0.4$, al fondo della crosta $\epsilon_b\approx\rho_n\approx m_nn_n$, $k_{F,n}=(3\pi^2n_n)\expy{\frac{1}{3}}$, densit\'a saturazione nucleare $n_0=\SI{0.15}{\per\cubic\femto\meter}$, ma nella materia nucleare $n_n\approx n_p\approx \frac{n_0}{2}$. La pressione \'e piccola rispetto alla densit\'a di energia perch\'e $E_F\ll m_m$. Per $n_n=n_0$ si ha $k_{F,n}\approx\SI{1.64}{\per\femto\meter}$, usando $\hbar c\approx\SI{197}{\mega\ev\femto\meter}$ si ha $k_{F,n}\approx\SI{324}{\mega\ev}$; a $\frac{n_0}{2}$ $K_{F,n}\approx\SI{257}{\mega\ev}$
                \end{column}
            \end{columns}
                   
        \item \xaumenta{M_*}, \xdiminuisce{\Delta R_{crust}} - Usando la relazione $\frac{\Delta R}{R}\approx \frac{1-2C}{C}\frac{P_b}{\epsilon_b}$: \SI{80}{\meter} for a $2\msun{}$ star with $R=\SI{10}{\kilo\meter}$ and $C=0.3$, and \SI{400}{\meter} for a $1.4\msun{}$ star with $R=\SI{15}{\kilo\meter}$ and $C=0.15$.
            
    \end{itemize}
\end{frame}

\begin{frame}{Pasta in the Crust}
    \begin{itemize}
        \item Pasta in the Crust (Ravehall '83)- Asymmetrical Wigner-Seitz cells: cells gettin closer can marge in one direction to form rods of dense nuclear matter (dnm) - several chains along same direction form spaghetti-like structure surrounded by liquid neutron. Increasing further density many rods form a plate, givin rise to a lasagna-like structure of dnm plates surrounded by pure neutron matter plate. A densit\'a maggiori si invertono i ruoli tra d.n.m. e dripping di liquido neutronico: lasagna \'e self-inverted, per gli spaghetti ho rods di neutroni liquidi circondati da dnm, infine bolle di neutroni liquidi circondati da dnm.
        \item Condizioni stabilit\'a pasta-phase: energia singola cella $\frac{E_{tot}}{V}=\epsilon_t=\epsilon_b+\epsilon_s+\epsilon_C+\epsilon_e$: bulk, superficie, Coulomb, elettroni, minimizzo risp. raggio $R_N$ della struttura nucleare: bulk e electron non dipendono da $R_N$ quindi non influenzano il minimo. $dq=4\pi r^2\rho\,dr$, $\phi(r)=\frac{1}{4\pi\epsilon_0}\frac{q(r)}{r}$, en. potenziale $dU=\phi(r)\,dq$: $U\propto \frac{Q^2}{R_N}$.
        \item $n_e(r)$ uniforme (lunghezza onda compton $(\frac{\hbar}{m_ec}=)\frac{1}{2\pi m_e}\approx\SI{e-12}{\meter}$, $R_{WS}\approx\SIrange{10}{20}{\femto\meter}$, o pi\'u formalmente $C_C^{screen}=\frac{e^2}{r}\exp{-\frac{r}{\lambda_{TF}}}$, dove $\lambda_{TF}=\frac{\pi^2}{2e}(k_F\sqrt{k_F^2+m_e^2})\expy{-\frac{1}{2}}$ e $k_F=(3\pi^2\rho_e)\expy{\frac{1}{3}}$: se $\lambda_{TF}\gg L$, lunghezza d'onda della materia nucleare, il gas elettronico non riesce a schermare localmente rods etc e resta quasi uniforme).
        \item $\epsilon_{surf}\propto \frac{R_N^2}{R_N^3}=\frac{1}{R_N}$, $\epsilon_C\propto \frac{R_N^5}{R_N^3}=R_N^2$.
        \item Minimiziamo la somma dei due termini: $\TDof{R_N}(\frac{C_s}{R_N}+C_cR_N^2)=0$ che da $\epsilon_s=2\epsilon_c$ e questo vale per le paste di dimensione superiore considerando la densit\'a di energia.
            \begin{columns}[T]
                \begin{column}{0.55\textwidth}
\begin{align*}
& Q \propto R_N^3,\ E_c\propto\frac{Q^2}{R_N}\propto R_N^5:\ \epsilon_c\propto R_N^2\tag{sfera}\\
& \frac{Q}{L} \propto R_N^2,\ \frac{E_c}{L}\propto\left(\frac{Q}{L}\right)^2\propto R_N^4:\ \epsilon_c\propto R_N^2\tag{rod}\\
& \frac{Q}{A} \propto R_N,\ \frac{E_{\rm C}}{A}\propto\left(\frac{Q}{A}\right)^2 R_N\propto R_N^3:\ \epsilon_c\propto R_N^2\tag{slab}
\end{align*}
                \end{column}
                \begin{column}{0.45\textwidth}
                    rod di raggio $R_N$: $S_{rod}=2\pi R_N^2L$, $V_{rod}=\pi R_N^2L$ quindi $\frac{S_{rod}}{V_{rod}}\propto \frac{1}{R_N}$, per slab $S_{slab}=A^2$, $V_{slab}=2R_NA^2$ quindi $\frac{S_{slab}}{V_{slab}}\propto \frac{1}{R_N}$.
                    Crosta interna ricca di neutroni ($N\gg Z$) bisogna sapere come la tensione superficiale cambia con l'asimmetria di isospin
                    
                \end{column}
            \end{columns}
    \end{itemize}
\end{frame}

\section{Neutron Star Core}

\subsection{Outer core}

\begin{frame}{Nuclear Matter}
    \begin{itemize}
        \item Infinite Symmetric nuclear matter: $A\to\infty$, $N=Z$, interazione coulombiana spenta. Binding energy per particles $\frac{E_b}{A}=a_v+a_sA\expy{-\frac{1}{3}}+a_cZ^2A\expy{-\frac{2}{3}}+a_{sym}\frac{(N-Z)^2}{A^2}+\frac{\delta(A)}{A}$ che si riduce a $\frac{E_b}{A}|_{n.m.}=a_V=-\SI{16}{\mega\ev}$.
        \item Esperimenti di scattering su $^{208}Pb$ mostrano che dopo la superficie di \SI{1}{\femto\meter} la densit\'a \'e costante $n_0=\frac{A}{V}=\frac{3A}{4\pi R^3}=\frac{3}{4\pi r_0^3}$: fit from mean-field model $n_0\approx\SIrange{0.15}{0.16}{\per\cubic\femto\meter}$ or $\rho_0=m_Nn_0\approx\SIrange{140}{150}{\mega\ev\per\cubic\femto\meter}$, ie $\rho_0\approx\SIrange{2.5e14}{2.7e14}{\gram\per\cubic\cm}$, where $m_N=\SI{939}{\mega\ev}$. The Fermi momentum at saturation density is $k_{F,0}=(\frac{6\pi^2n_0}{\gamma})\expy{\frac{1}{3}}=\SI{260}{\mega\ev}$. $\gamma$ \'e il fattore di degenerazione per neutroni e protoni: $2I+1$ isospin e $2S+1$ di spin.
        \item EOS of nuclear matter $\epsilon(n_B)$: $\frac{E_b}{A}(n_B)=\frac{\epsilon(n_B)}{n_B}-m_N$, we also know that $\frac{\epsilon(n_0)}{A}=-\SI{16}{\mega\ev}$. As nuclear matter is stable at $n_0$: $\PDy{n_0}{\epsilon}|_{n_0}=0$ (energia per nucleone $e(n_B)=\frac{\epsilon(n_B)}{n_B}$, quindi $\epsilon=n_Be$, quindi $\TDy{n_B}{\epsilon}=e+n_N\TDy{n_B}{e}$ e $P=n_B(e+n_B\TDy{n_B}{e})-n_Be$ da cui $P=n_B^2\TDy{(E/A)}{n_B}$).
        \item Incompressibility - $K_{\infty}=k_{F,0}^2\TtwoDy{k_{F,0}}{(E/A)}|_{k_{F,0}}=9n_0^2\TtwoDy{n_0}{(E/A)}|_{n_0}$: related to excitation energy of nuclei (giant monopole, dipole res, $E_{ex}$ several \si{\mega\ev}, compression mode). Il valore sperimentale ricavato da fit della leptodermus expansion $K_A\approx K_{\infty}+K_{surf}A\expy{-\frac{1}{3}}+K_{\tau}(\frac{N-Z}{A})^2+K_C \frac{Z^2}{A\expy{\frac{4}{3}}}+\ldots$ con dati risonanze da $K_{\infty}\approx\SI{240+-20}{\mega\ev}$. $K_{\infty}$ non controlla espansione EOS in potenze di $(n-n_0)$ per alte densit\'a, ma solo attorno a densit\'a di saturazione $n_0$.
        \item Hugenholtz-van Hove theorem - Fermi energy of Fermi gas with interactions at $T=0$ \'e l'energia media pi\'u la pressione per particella: $E_F=\TDy{n_B}{\epsilon}=\frac{\epsilon+P}{n_B}$; using $\epsilon=-P+\mu n$ so at $T=0$ Fermi energy is equal to chemical potential even if interactions are included.
        \item At saturation density $P=0$: $E_F(n_0)=\frac{\epsilon}{n_B}|_{n_0}=\frac{\epsilon_b}{A}(n_0)+m_N=\SI{923}{\mega\ev}$
    \end{itemize}
\end{frame}

\begin{frame}{Dirac phenomenology: effective mass to include interactions}
    \begin{columns}[T]
        \begin{column}{0.4\textwidth}
            \begin{figure}[!ht]\includegraphics[trim={0cm 0cm 0 0},clip,keepaspectratio,width=0.95\textwidth]{/Schaffner/EOS_kinf_mstar}\label{fig:/Schaffner/EOS_kinf_mstar}\end{figure}
        \end{column}
        \begin{column}{0.6\textwidth}
    \begin{itemize}
        \item Dirac phenomenology - $E_F=\sqrt{k_F^2+m^2}$, which will be modified when including interactions: shift in the mass of particle by some scalar potential $m^*=m-S(n)$ and 4-moment shifted by vectorial potential $\tensor{k^*}{^\mu}=\tensor{k}{^\mu}-\tensor{V}{^\mu}(n)$.
        \item Consideriamo shift di energia con materia nucleare isotropa (spatial components of V null): $E_F=\sqrt{k_F^2+(m+S(n))^2}+\tensor{V}{^0}(n)=\sqrt{k_F^2+m^{*2}}+\tensor{V}{^0}(n)$ - at low density scalar potential dominate so nuclear matter is bound, at high density the repulsive vector potential becomes larger than scalar so nuclear matter get unbound, ie pressione positiva alla compressione: at saturation density $n_0$ $E_{F,0}=\SI{923}{\mega\ev}$ so if we fix $m^*$ we fix $\tensor{V}{^0}$.
    \end{itemize}
            
        \end{column}
    \end{columns}
    \begin{itemize}
        \item Combined potential felt by nucleons within a nucleus at zero momentum is $U_N(n_0)=S(n_0)+\tensor{V}{^0}(n_0)\approx-\SI{50}{\mega\ev}$; infatti usando il T. HvH, in NR limit, $\frac{E_b}{A}=E_{kin}(n_0)+U_N(n_0)=\frac{k_{F,0}^2}{2m}+U_N(n_0)$ ovvero $U_N(n_0)\approx \frac{E_b}{A}-\frac{k_{F,0}}{2m}=(-16-35)\si{\mega\ev}\approx-\SI{50}{\mega\ev}$ con $k_{F,0}\approx\SI{260}{\mega\ev}$.
    \end{itemize}
\end{frame}

\begin{frame}{Pure neutron matter}
    \begin{columns}[T]
        \begin{column}{0.45\textwidth}
    \begin{itemize}
        \item Frazione protonica $x_p=\frac{Z}{A}$: per materia nucleare simmetrica $x_p=\frac{1}{2}$, per pure neutron matter $x_p=0$
        \item Asimmetria $\delta=\frac{n_n-n_p}{n_n+n_p}=1-2x_p$: $\delta=1$ pure neutron matter, $\delta=0$ materia simmetrica
        \item Neutron Drip line alla fine della crosta esterna: $\frac{Z}{A}\approx0.3$ quindi $\delta=1-2*0.3=0.4$; fine crosta interna cella WS ha $Z=40$, $N=1800$: $\frac{Z}{A}\approx\num{0.022}$, $x_p\approx\num{0.02}$, $\delta\approx\num{0.956}$
        \item Nuclei N-rich conosciuti: $^8He$ $\frac{Z}{A}=0.25$,$^{11}Li$ $\frac{Z}{A}\approx0.27$, $^{14}Be$ $\frac{Z}{A}\approx0.29$ sono ancora lontani da $\frac{Z}{A}\approx\num{0.02}$ - I nuclei N rich sono esotici
        \item I nuclei doppiamente magici usati nel fit della formula semiempirica di massa: $^{132}Sn$, $\frac{50}{132}\approx\num{0.38}$, $^{208}Pb$, $\frac{82}{208}\approx\num{0.39}$
        \item I nuclei ordinari vincolano bene a bassi valori di $\delta$: $\frac{E_b}{A}(n,\delta)\approx \frac{E}{A}(n,0)+S(n)\delta^2+\ldots$ mentre PNM corrisp. $\delta=1$
    \end{itemize}
        $\frac{E_{PNM}}{A}-\frac{E_{SNM}}{A}\approx J$: differenza tra $\delta=0$ e $\delta=1$
            \begin{align*}
                &\frac{E_b}{A}(n,\delta)\approx E_0(n)+S(n)\delta^2,\ E_0(n_0)\approx-\SI{16}{\mega\ev}\\
                &\frac{E}{A}(n_0,\delta=1)\approx-\SI{16}{\mega\ev}+J\Rightarrow \frac{E}{A}|_{PNM}(n_0)\approx\SI{16}{\mega\ev}
            \end{align*}
        \end{column}
        \begin{column}{0.55\textwidth}
            \begin{itemize}
                \item Simmetria forza forte per $\Pproton\leftrightarrow\Pneutron$, o stabilit\'a materia simmetrica: $\frac{E}{A}(n_B,\delta)=\frac{E}{A}(n_B,0)+\frac{1}{2}\PtwoDy{\delta}{E/A}|_{\delta=0}\delta^2+O(\delta^4)$. Symmetry energy $S(n_B)=\frac{1}{2}\PtwoDy{\delta}{E/A}|_{\delta=0}$
                \item Espansione attorno $n_0$:
                    \begin{align*}
                        &S(n_B)\approx S(n_0)+\PDy{n_B}{S}|_{n_0}(n_B-n_0)+\ldots\\
                        &=J+L \frac{n_B-n_0}{3n_0}+\ldots\\
                        &J=S(n_0)\approx\SI{32.3+-0.5}{\mega\ev}\tag{sym energy sat inf nuc mat}\\
                        &L=3n_0\PDy{n_0}{S}|_{n_0}\approx\SI{53.5+-15.0}{\mega\ev}\tag{pendenza risp densit\'a a $n_0$}
                    \end{align*}
                \item $P_{PNM}(n_0)\approx \frac{n_0L}{3}$
                    %Per materia simmetrica: $P=n^2\PDy{n}{E/A}$ ($E_0=\frac{E}{A}(n,\delta=0)$) quindi $P_{SNM}(n_0)=0$. Per materia asimmetrica: $\frac{E}{A}(n,\delta)\approx E_0(n)+S(n)\delta^2$ quindi $\PDy{n}{E/A}=E_0'(n)+S'(n)\delta^2$, $E_0'(n_0)=0$, $S?(n_0)\neq0$ quindi $P(n_0,\delta)\approx n_0^2S'(n_0)\delta^2$
                    \begin{align*}
                        &E_0=\frac{E}{A}(n,\delta=0)\tag{materia simmetrica}\\
                        &P=n^2\PDy{n}{(E/A)}\ \Rightarrow\ P_{SNM}(n_0)=0\\
                        &\frac{E}{A}(n,\delta)\approx E_0(n)+S(n)\delta^2\tag{materia asimmetrica}\\
                        &\PDy{n}{E/A}=E_0'(n)+S'(n)\delta^2\Rightarrow P(n_0,\delta)\approx n_0^2S'(n_0)\delta^2\\
                        &S'(n_0)=\frac{L}{3n_0}:\ P(n_0,\delta)\approx \frac{n_0L}{3}\delta^2\\
                        &a_{sym}\approx\SI{23}{\mega\ev}<J\tag{SM: sym term has bulk+surf}
                    \end{align*}
            \end{itemize}
        \end{column}
    \end{columns}
\end{frame}

\begin{frame}{$L$ determina pressione PNM vicino $n_0$}
    \begin{itemize}
        \item $P=n_B^2\PDy{n_B}{(E/A)}$: $\frac{E}{A}|_{PNM}\approx E_0(n)+S(n)$, alla saturazione $E'_0(n_0)=0$, quindi $P_{PNM}=n_0^2S'(n_0)$ ma $S'(n_0)=\frac{L}{3n_0}$: $P_{PNM}(n_0)=\frac{L}{3}n_0$.
        \item Confronto tra pressione PNM vicino $n_0$ e densit\'a di energia: $\frac{\epsilon}{n_0}=m_N+\frac{E_b}{A}\approx m_N-\SI{16}{\mega\ev}+J$ quindi $\epsilon\approx n_0(\SI{923}{\mega\ev}+J)$, il rapporto $\frac{P}{\epsilon}=\frac{L/3n_0}{(\SI{923}{\mega\ev}+J)n_0}$. Con $L=\SI{53.5}{\mega\ev}$, $J=\SI{23.3}{\mega\ev}$ si ha $\frac{P}{\epsilon}\approx\num{0.019}$: la pressione di PNM \'e molto minore di densit\'a di energia a $n_0$.
        \item With ratio $P/\epsilon$ we can estimate thickness of outer part of NS up to $n_0$: using $\frac{\Delta R}{R}\approx \frac{1-2C}{C}\frac{P}{\epsilon}\approx\numrange{3}{6}\%$, con $C=\numrange{0.2}{0.3}$: most of NS is beyond saturation. Pressure is very low at $n_0$: increase of P and stiffening of EOS is needed for NS of $2\msun$.
        \item Quali NS arrivano fino a $n_0$: $\frac{\Delta R}{R}\approx1$ e usando espressione precedente si trova $C\approx \frac{1}{50}\approx0.02$. Usnado modello pi\'u accurato: $M\approx0.2\msun$, $R\approx\SI{20}{\kilo\meter}$, la compattezza $C=\frac{GM}{c^2R}$, usando $\frac{G\msun}{c^2}\approx\SI{1.477}{\kilo\meter}$, quindi $C\approx \frac{1}{70}$ vicino a stima grezza. La fisica a $n_0$ non basta a sostenere normale NS.
    \end{itemize}
\end{frame}

\begin{frame}{Scattering elastico}
    \begin{columns}[T]
        \begin{column}{0.5\textwidth}
            \begin{align*}
                &\psi(\vec{r})\xrightarrow{r\to\infty}\exp{ikz}+f(\theta)\frac{\exp{ikr}}{r}=\psi_{inc}+\psi_{sc}\tag{sc wf asint}\\
                &f(\theta)=\sum_{l=0}(2l+1)f_l(k)P_l(\cos{\theta})\tag{sviluppo onde parziali}\\
                &f_l(k)=\frac{\exp{2i\delta_l}-1}{2ik}=\frac{\exp{i\delta_l}\sin{\delta_l}}{k}=\frac{1}{k\cot{\delta_l}-ik}\\
                &k\cot{\delta_0(k)}=-\frac{1}{a_s}+\frac{1}{2}r_ek^2+\ldots\tag{effective range expansion}\\
                &\Rightarrow f_0(k)\approx \frac{1}{-\frac{1}{a_s}+\frac{1}{2}r_ek^2-ik}\xrightarrow{k\to0}-a_s
            \end{align*}
        \end{column}
        \begin{column}{0.5\textwidth}
            \begin{itemize}
                \item Poli $f_0(k)$: $-\frac{1}{a}-ik=0$ quindi $k=\frac{i}{a}$ - per $a>0$ la funzione d'onda radiale asintotica $\propto\exp{-\kappa r}$, con $\kappa=\frac{1}{a}$, funzione d'onda normalizzabile - stato legato $E=-\frac{\hbar^2\kappa^2}{2\mu}<0$; per $a<0$ $k=i\kappa$ con $\kappa=\frac{1}{|a|}$ - stato virtuale con funzione d'onda non normalizzabile ma $E<0$ (virtual state \'e un tipo specifico di polo sottosoglia).
            \end{itemize}
        \end{column}
    \end{columns}
    \begin{itemize}
        \item In s-wave $f(\theta)\approx f_0(k)$: sezione d'urto differenziale isotropa $\TDy{\Omega}{\sigma}=|f_0|^2$ - $\sigma_s=\int\,d\Omega|f_0|^2=4\pi|f_0|^2$; $f_0=\frac{\exp{i\delta_0}\sin{\delta_0}}{k}$ quindi $\sigma_s=\frac{4\pi}{k^2}\sin^2{\delta_0}$ e $|f_0|_{max}=\frac{1}{k}$ da cui $\sigma_a^{max}=4\pi \frac{1}{k^2}$.
    \end{itemize}
\end{frame}

\begin{frame}{Pure Nuclear Matter EOS: from 2B to bulk matter}
            \begin{itemize}
                \item A basse densit\'a le propriet\'a di PNM possono essere ricavate da scatterin a due corpi NN:
                \item I nucleoni sono fermioni di spin $\frac{1}{2}$ quindi la funzione d'onda del sistema di due nucleoni \'e antisimmetrica sotto scambio di coppie: $P_{12}\psi=-\psi$. $\Psi=\psi(\vec{r}_1,\vec{r}_2)\chi_{spin}\chi_{iso}$: stato s-wave $L=0$ quindi $\chi_{spin}\chi_{iso}$ antisimmetrica, $S=1$, $I=0$ ($^3S_1$ spin triplet) o $S=0$, $I=1$ ($^1S_0$ spin singlet); per nn si ha necessariamente $I=1$, $I_3=-1$ quindi isospin simmetrico e spin antisimmetrico; il canale $^3S_1$, $I=0$ \'e il deutone \Pproton\Pneutron (pi\'u componente $^3D_1$ dovuto alla forza tensoriale)
                \item \Pneutron\Pneutron non forma stato legato ma c'\'e ''risonanza sotto-soglia''/near-threshold virtual state: $a_{nn}\approx-\SI{18.5}{\femto\meter}$, $r_e\approx\SI{2.7}{\femto\meter}$, $|a_{nn}|\gg r_e$ cio\'e scattering length maggiore del range di interazione nucleare quindi il sistema \Pneutron\Pneutron nel canale $^1S_0$ \'e vicino alla formazione di uno stato legato, $k_{pole}=-\frac{i}{|a|}$ stato virtuale vicino a $E=0$ (distanza energetica caratteristica dalla soglio $\frac{\hbar^2}{m_na^2}\approx\SI{0.1}{\mega\ev}$): $a<0$ indica che non \'e presente un vero shallow bound state - se $a_s$ fosse grande e positivo avrei stato debolmente legato con $E_B=\frac{\hbar^2}{ma^2}$
                \item Unitary gas - regime quasi-universale del gas di fermi fortemente interagente: si ha quando $k_F|a_s|$ grande e $k_Fr_e$ rimane relativamente piccolo e $r_e\ll\invers{k_F}$ cio\'e distanze molto maggiori di range forza forte. $a\gg \frac{1}{k_F}$: scattering length molto pi\'u grande della distanza caratteristica associata alla densit\'a, $f_0(k)\approx \frac{1}{-ik}=\frac{i}{k}$ - lo scattering non dipende pi\'u dai dettagli microscopici del potenziale nucleare, $\sigma_s=4\pi|f_0|^2\approx \frac{4\pi}{k^2}$ (limite massimo consentito da unitariet\'a per s-wave).
                \item Densit\'a critica $a_sk_F\approx1$, regime quasi-unitario $k_Fa_s\gg1$: $n=\frac{k_F^3}{3\pi^2}$, $d\approx n\expy{-\frac{1}{3}}$, quindi $d\approx \frac{3.09}{k_F}\approx \frac{\pi}{k_F}$ - $k_F\approx \frac{\pi}{a_s}\approx\SI{33.5}{\mega\ev}\approx\SI{0.17}{\per\femto\meter}$ (pi\'u precisamente $d=\frac{\lambda_F}{2}=\frac{\pi}{k_F}$).
                \item Range densit\'a unitary Fermi gas in NS: $\num{e-3}n_0(\text{ordine grandezza dripline})<n<0.1n_0$ - A densit\'a pi\'u alte non basta 2B interaction
            \end{itemize}
\end{frame}

\begin{frame}{Unitary Fermi Gas: gas diluito fortemente interagente e oltre}
    \begin{columns}[T]
        \begin{column}{0.6\textwidth}
    \begin{itemize}
        \item Biblio: ''Low-Density Neutron Matter and the Unitary Limit''
        \item Quando $a\to\infty$, $r_e\to0$ non rimane alcuna lunghezza microscopica dell'interazione: la sola scala di lunghezza del sistema a molti corpi \'e $\invers{k_F}$
        \item Energia proporzionale a energia di Fermi $\frac{E}{N}\propto E_F$: uni scala disponibile $E_F=\frac{\hbar^2k^2}{2m_n}$ (gas Fermi libero a $T=0$ $(\frac{E}{N})_{FG}=\frac{3}{5}E_F$): $(\frac{E}{N})_{UG}=\xi(\frac{E}{N})_{FG}=\xi \frac{3\hbar^2k_F^2}{10m_n}$, $\xi$ \'e il Bertsch parameter.
        \item \(\xi\) non dipende dal dettaglio del potenziale microscopico: nel limite unitario \'e un numero universale per gas di fermioni a due componenti con interazione a corto raggio. Gli esperimenti con atomi ultrafreddi, ie \(^{6}\mathrm{Li}\), possono portare il sistema molto vicino a $ |a|\to\infty$ mediante una Feshbach resonance: $\xi\approx\num{0.37}$.
        \item Neutron Matter reale: correzioni $\invers{a_{nn}}$, $r_e$, etc: $\frac{E}{N}=\frac{3}{5}E_F[\xi+\frac{c_1}{k_Fa_{nn}}+c_2k_Fr_e+\ldots]$
        \item Approx Unitary Fermi Gas valida per $n\leq\SI{0.02}{\per\femto\meter}\approx\frac{0.02}{0.16}n_0$ circa $10\%n_0$: non bastano $a_s$ e $r_e$ ma servono anche interazioni a 3-corpi calibrate su $^3H$, $^3He$, $^4He$
        \item Chiral effective field theory $\chi$EFT: si tengono gradi di libert\'a rilevanti a scala energetica considerata (EFT) $V=V_{LO}+V_{NLO}+V_{N^2LO}+\ldots$ e la simmetria chirale che vincola la forma delle interazioni nucleoni-pioni.
    \end{itemize}
            
        \end{column}
        \begin{column}{0.4\textwidth}
            \begin{figure}[!ht]\includegraphics[trim={0cm 0cm 0 0},clip,keepaspectratio,width=0.95\textwidth]{Schaffner/EOS_chiEFT}\label{fig:Schaffner_EOS_chiEFT}\end{figure}
            \begin{itemize}
                \item Fino $n\approx\SI{0.2}{\per\cubic\femto\meter}$ $E/A$ sono dello stesso ordine di grandezza
                \item La curva inferiore \'e unitary Fermi gas: salendo in densit\'a i calcoli realistici includono effetti del range finito, partial waves superiori e effetti a 3 neutroni.
            \end{itemize}
        \end{column}
    \end{columns}
    
\end{frame}

\begin{frame}{Collegamento symmetric nuclear matter to pure nuclear matter}
\begin{columns}[T]
    \begin{column}{0.5\textwidth}
    \begin{itemize}
        \item Ansatz $\frac{E_b}{A}(n_B,\delta)=\frac{E_b}{A}(n_B,0)+S(n_B)\delta^2$, come suggerito da simmetria interazione forte per $\Pproton\leftrightarrow\Pneutron$ $E(\delta)=E(-\delta)$ \'e funzione pari.
    \end{itemize}
        
    \end{column}
    \begin{column}{0.5\textwidth}
    \begin{table}
\centering
\begin{tabular}{lcc}
\hline
\textbf{Modello / approccio}
& \(\boldsymbol{J}\) [MeV]
& \(\boldsymbol{L}\) [MeV]
\\
\hline
FRDM2012
& \(32.3 \pm 0.5\)
& \(53.5 \pm 15\)
\\
Chiral EFT
& \(28.9\text{--}34.9\)
& \(43.0\text{--}66.0\)
\\
\hline
\end{tabular}
\end{table}
        
    \end{column}
\end{columns}
    \begin{itemize}
        \item Ponendo $\delta=1$ si ha PNM: $\frac{E_b}{A}(n_B,1)=\frac{E_b}{A}(n_N,0)+S(n_B)$ quindi in approssimazione quadratica $S(n_B)\approx \frac{E}{A}|_{PNM}-\frac{E}{A}|_{SNM}$ - vincoliamo la symmetry energy con calcoli microscopici per PNM.
        \item Alla saturazione $J=S(n_0)\approx\frac{E}{A}|_{PNM,n_0}-\frac{E}{A}|_{SNM,n_0}$, poich\'e $\frac{E}{A}|_{SNM,n_0}\approx-\SI{16}{\mega\ev}$ e $\frac{E}{A}\approx\SIrange{15}{20}{\mega\ev}$: $J\approx\SIrange{30}{35}{\mega\ev}$; $L\approx3n_0\TDy{n}{(E/A)_{PNM}}|_{n_0}$.
    \end{itemize}

\end{frame}

\subsection{Relazione massa-raggio e stabilit\'a}

\begin{frame}{Densit\'a di energia all'aumentare della densit\'a }
    \begin{columns}[T]
        \begin{column}{0.6\textwidth}
            \begin{itemize}
                \item Soluzione TOV (equilibrio idrostatico$\neq$equilibrio dinamico) risolvendo dal centro con $P_c=P(\epsilon_c)$.
                \item Transition OC-IC: $\frac{\epsilon}{\epsilon_0}\approx0.002$ different model give almost constant value, Crust-Core: $\frac{\epsilon_c}{\epsilon_0}\approx0.3$ differs from model to model - Transition point is not well determined but EOS is almost the same, ie PNM, for all models.
                \item EOS of Outer Crust is power-law, the onset of inner crust show sudden decrease of slope due onset of liquid neutrons. In Core pressure increases strongly with energy density.
                \item A: WD reach maximum mass for $\frac{\epsilon_c}{\epsilon_0}\approx\num{4e-6}$ (\SI{e9}{\gram\per\cubic\cm}), D-E: stable NS, $\frac{\epsilon_c}{\epsilon_0}\approx\numrange{0.5}{8}$ (EOS known up to \num{1.3}).
                \item $\epsilon=\epsilon_{rest}+\epsilon_{kin}+\epsilon_{int}$: baryon number density start to deviate from $\rho\approx\epsilon$ at $n_0$ and as interaction between neutrons has to be repulsive the increase in number density is lower than that in energy density.
            \end{itemize}
        \end{column}
        \begin{column}{0.4\textwidth}
            \begin{figure}[!ht]\includegraphics[trim={0cm 0cm 0 0},clip,keepaspectratio,width=0.95\textwidth]{Schaffner/EOS_WD_NS}\label{fig:Schaffner_EOS_WD_NS}\end{figure}
            \begin{figure}[!ht]
                \includegraphics[trim={0cm 0cm 0 0},clip,keepaspectratio,width=0.95\textwidth]{Schaffner/energydensity_vs_baryon}\label{fig:Schaffner_energydensity_vs_baryon}
            \end{figure}
        \end{column}
    \end{columns}
    
\end{frame}

\begin{frame}{Polytropic power (adiabatic index) vs energy density}
    \begin{columns}[T]
        \begin{column}{0.6\textwidth}
            \begin{itemize}
                \item Una transizione di fase di primo ordine \'e caratterizzata dalla discontinuità di una derivata prima del potenziale termodinamico. A \(T=0\), per due fasi in equilibrio vale $P_1=P_2$, $\mu_1=\mu_2$, mentre la densit\'a barionica pu\'o presentare un salto, $n_1\neq n_2$, e analogamente pu\'o essere discontinua la densit\'a di energia. Nella regione di coesistenza idealizzata la pressione resta costante mentre cambia la densit\'a, per cui $\Gamma=\frac{n_B}{P}\frac{dP}{dn_B}\simeq 0$.
            \item Polytropic power $\Gamma=\TDly{n_B}{P}$: transizione del prim'ordine, cambia densit\'a a pressione costante, quindi $\Gamma\approx0$.
            \item From the left: dip for transition $^{56}Fe\to^{62}Ni\to^{64}Ni$ in outer crust; after A we have transition to heavier and more neutron rich nuclei with deeper dip.
            \end{itemize}
        \end{column}
        \begin{column}{0.4\textwidth}
            \begin{figure}[!ht]\includegraphics[trim={0cm 0cm 0 0},clip,keepaspectratio,width=0.95\textwidth]{Schaffner/Gamma_vs_energydensity}\label{fig:Schaffner_Gamma_vs_energydensity}
            \end{figure}
        \end{column}
    \end{columns}
    \begin{itemize}
        \item Even if there is no dip $\Gamma\approx \frac{4}{3}$, lower limit for stability of polytropes, at beginnin of inner crust $\Gamma\approx0.5$ due to onset of neutron liquid: when neutron level starts to be populated occupy lowest levels in  momentum first, giving no pressure contrib. but n_B keeps increasing so that $\Gamma$ drops.
        \item The wiggle before the core is due to transition from lattice structure of the crust to uniform liquid in the core; max of $\Gamma\approx3.5$ is reached for density of core of NS; then drops to two showing EOS dominated by interactions.
        \item The correct EOS for high density is not the one based on nucleons but rather the EOS taking into account of quarks and gluons: the correct high density limit $\Gamma=\frac{4}{3}$ is the value of UR free particles.
    \end{itemize}
\end{frame}

\begin{frame}{Stabilit\'a dinamica nel diagramma Massa-Raggio}
            \begin{figure}[!ht]
                \includegraphics[trim={0cm 0cm 0 0},clip,keepaspectratio,width=0.3\textwidth]{Schaffner/M_vs_energydensity}\label{fig:Schaffner_M_vs_energydensity}~
                \includegraphics[trim={0cm 0cm 0 0},clip,keepaspectratio,width=0.3\textwidth]{Schaffner/R_vs_energydensity}\label{fig:Schaffner_R_vs_energydensity}~
                \includegraphics[trim={0cm 0cm 0 0},clip,keepaspectratio,width=0.3\textwidth]{Schaffner/M_vs_R}\label{fig:Schaffner_M_vs_R}
            \end{figure}
            \begin{itemize}
                \item Stab. sol. eq. idro. a perturbazioni radiali $\xi(r,t)=\xi_n(r)\exp{i\omega_nt}$: eq. di E. linearizzate vincolano perturbazioni metriche \(\delta\Phi,\delta\Lambda\) e della massa \(\delta m\) alle perturbazioni della materia; conservazione barionica e EOS permettono di esprimere \(\delta P\) e \(\delta\varepsilon\) mediante spostamento radiale \(\xi\). Sostituendo queste relazioni in eq. radiale del moto, rimane una equazione differenziale per \(\xi(r,t)\): decomposizione in modi normali da Sturm--Liouville.
                \item Autovalori $\omega_0^2<\omega_1^2<\omega_2^2<\ldots$, autof. modo n con autov. $\omega_n$ possiede n nodi radiali: $\omega_n^2>0$ modo oscillante, $\omega_n<0$ modo radialmente instabile.
                \item Per EOS fissata si costruisce famiglia di NS al variare della densit\'a di energia centrale $\epsilon_c$: risolvendo TOV si ottiene $M(\epsilon_c)$, $R(\epsilon_c)$. Turning-point argument: cambia numero di modi instabili quando M ha estremo $dM/d\epsilon_c=0$. Su un ramo normalmente stabile $dM/d\epsilon_c>0$, regione stabile, e attraversando un estremo di M un $\omega_n^2$ cambia segno: se $dR/d\epsilon_c|_{extr}<0$ cambia segno un modo pari, altrimenti un modo dispari. A:Limite Chandra - $\omega_0^2(<\omega_1^2)>0\ldots$, $M(\epsilon_c)$ raggiunge max($dM/d\epsilon_c=0$), $dR/d\epsilon_c<0$ cambia segno modo pari, ie $\omega_0^2$, quindi si entra in ramo instabile; B: M ha minimo con $dR/d\epsilon_c>0$, quindi cambia segno modo dispari, ie $\omega_1^2<0$; C: max in M while R incresing with $\epsilon_c$ quindi cambia segno modo dispari, ie $\omega_1^2>0$; D: ramo NS, M ha min e $dR/d\epsilon_c<0$, cambia segno modo pari, ie $\omega_0^2>0$, $M_{min}=0.09\msun$, $R=\SI{250}{\kilo\meter}$; E: massa max NS, $dR/d\epsilon_c<0$ quindi cambia segno modo pari, ie $\omega_0^2<0$, entriamo in ramo instabile - $M=2.4\msun{}$ max NS.
            \end{itemize}
    
\end{frame}

\begin{frame}{Direttamente da curva $M(R)$}
            \begin{itemize}
                \item Si percorre la curva aumentando la densit\'a centrale da destra a sinistra
                \item Percorrenza antioraria \'e equivalente a raggio decrescento all'aumentare di $\epsilon_c$: antioraria-modo diventa instabile, oraria-modo diventa stabile
            \end{itemize}
            \begin{columns}[T]
                \begin{column}{0.5\textwidth}
            \begin{figure}[!ht]
                \includegraphics[trim={0cm 0cm 0 0},clip,keepaspectratio,width=0.95\textwidth]{Schaffner/M_vs_R}\label{fig:Schaffner_M_vs_R2}
            \end{figure}
                \end{column}
                \begin{column}{0.5\textwidth}
            \begin{itemize}
                \item A: max M, antiorario - modo diventa instabile. Sequenza dopo A \'e instabile.
                \item B: min M, antiorario - modo diventa instabile
                \item C: max M, orario - modo diventa stabile
                \item D: min M; orario - modo diventa stabile. Here radius decreases with increasing central energy density. La sequenza dopo D \'e stabile.
            \end{itemize}
                \end{column}
            \end{columns}
\end{frame}

\subsection{Recap particelle}

\begin{frame}{Leptoni}
\begin{table}
\centering
\scriptsize
\resizebox{\textwidth}{!}{%
\begin{tabular}{lccclll}
\hline
\textbf{Leptone}
& \(\boldsymbol{mc^2}\) [MeV]
& \(\boldsymbol{Q}\)
& \textbf{Potenziale chimico}
& \(\boldsymbol{\tau_0}\) \textbf{(vuoto)}
& \textbf{Decadimenti principali}
& \textbf{Materia fredda}\\
\hline
\(e^-\)
& \(0.511\)
& \(-1\)
& \(\mu_e\)
& Stabile
& ---
& Presente, degenere\\

\(\mu^-\)
& \(105.66\)
& \(-1\)
& \(\mu_\mu=\mu_e\)
& \(2.197\times10^{-6}\,\mathrm{s}\)
& \(\mu^-\to e^-+\bar\nu_e+\nu_\mu\)
& Presente se \(\mu_e\geq m_\mu c^2\)\\

\(\tau^-\)
& \(1776.9\)
& \(-1\)
& \(\mu_\tau=\mu_e\)
& \(2.903\times10^{-13}\,\mathrm{s}\)
& \shortstack[l]{
  \(\tau^-\to\pi^-\pi^0\nu_\tau\) (\(25.5\%\))\\
  \(\tau^-\to e^-\bar\nu_e\nu_\tau\) (\(17.8\%\))\\
  \(\tau^-\to\mu^-\bar\nu_\mu\nu_\tau\) (\(17.4\%\))\\
  \(\tau^-\to\pi^-\nu_\tau\) (\(10.8\%\))
  }
& Generalmente assente\\
\hline
\(\nu_e\)
& ---
& \(0\)
& \(0\)
& ---
& Nessun decadimento osservato
& Non intrappolato\\

\(\nu_\mu\)
& ---
& \(0\)
& \(0\)
& ---
& Nessun decadimento osservato
& Non intrappolato\\

\(\nu_\tau\)
& ---
& \(0\)
& \(0\)
& ---
& Nessun decadimento osservato
& Non intrappolato\\
\hline
\end{tabular}%
}
\caption{Proprietà dei leptoni e principali decadimenti
nel vuoto. Tutti hanno spin \(1/2\) e numero barionico
\(B=0\). I neutrini sono considerati non intrappolati.}
\end{table}
\end{frame}

\begin{frame}{Mesoni}
    \begin{table}
\centering
\scriptsize
\resizebox{\textwidth}{!}{%
\begin{tabular}{lccccll}
\hline
\textbf{Mesone}
& \(\boldsymbol{mc^2}\) [MeV]
& \(\boldsymbol{Q}\)
& \(\boldsymbol{S}\)
& \textbf{Potenziale chimico}
& \(\boldsymbol{\tau_0}\) \textbf{(vuoto)}
& \textbf{Decadimenti principali (BR)}\\
\hline
\(\pi^-\)
& \(139.57\)
& \(-1\)
& \(0\)
& \(\mu_{\pi^-}=\mu_e\)
& \(2.603\times10^{-8}\,\mathrm{s}\)
& \(\pi^-\to\mu^-+\bar\nu_\mu\) (\(99.99\%\))\\

\(\pi^0\)
& \(134.98\)
& \(0\)
& \(0\)
& \(0\)
& \(8.43\times10^{-17}\,\mathrm{s}\)
& \shortstack[l]{
  \(\pi^0\to\gamma+\gamma\) (\(98.82\%\))\\
  \(\pi^0\to e^+e^-\gamma\) (\(1.17\%\))
  }\\
\hline
\(K^-\)
& \(493.68\)
& \(-1\)
& \(-1\)
& \(\mu_{K^-}=\mu_e\)
& \(1.238\times10^{-8}\,\mathrm{s}\)
& \shortstack[l]{
  \(K^-\to\mu^-+\bar\nu_\mu\) (\(63.56\%\))\\
  \(K^-\to\pi^-+\pi^0\) (\(20.67\%\))\\
  \(K^-\to\pi^-\pi^-\pi^+\) (\(5.58\%\))
  }\\

\(K^0,\bar K^0\)
& \(497.61\)
& \(0\)
& \(\pm1\)
& \(0\)
& \textit{Vedi sotto}
& \textit{Mixing } \(K^0-\bar K^0\)\\

\(K_S^0\)
& \(497.61\)
& \(0\)
& ---
& \(0\)
& \(8.954\times10^{-11}\,\mathrm{s}\)
& \shortstack[l]{
  \(K_S^0\to\pi^+\pi^-\) (\(69.2\%\))\\
  \(K_S^0\to\pi^0\pi^0\) (\(30.7\%\))
  }\\

\(K_L^0\)
& \(497.61\)
& \(0\)
& ---
& \(0\)
& \(5.1\times10^{-8}\,\mathrm{s}\)
& \shortstack[l]{
  \(K_L^0\to\pi^\pm e^\mp\nu\) (\(40.55\%\))\\
  \(K_L^0\to\pi^\pm\mu^\mp\nu\) (\(27.04\%\))\\
  \(K_L^0\to3\pi^0\) (\(19.52\%\))
  }\\
\hline
\end{tabular}%
}
\caption{Proprietà dei principali mesoni leggeri e decadimenti
nel vuoto. Tutti hanno \(B=0\), spin \(J=0\) e parità
\(P=-1\). Per i mesoni carichi positivi valgono
i canali di decadimento coniugati.}
\end{table}
I mesoni carichi positivi sono le antiparticelle
dei corrispondenti mesoni negativi e decadono
attraverso canali coniugati per carica
(ad esempio, $\pi^+\to\mu^++\nu_\mu$).
\end{frame}

\begin{frame}{Iperoni: propriet\'a e potenziale chimico in ColdNuclearMatter}
    \begin{table}
\centering
\scriptsize
\resizebox{\textwidth}{!}{%
\begin{tabular}{lccccccl}
\hline
\textbf{Iperone}
& \textbf{Quark}
& \(\boldsymbol{mc^2}\) [MeV]
& \(\boldsymbol{Q}\)
& \(\boldsymbol{S}\)
& \(\boldsymbol{J^P}\)
& \textbf{Potenziale chimico}
& \(\boldsymbol{\tau_0}\) \textbf{(vuoto)}
\\
\hline
\(\Lambda^0\) & \(uds\) & \(1115.68\) & \(0\)  & \(-1\)
& \(\frac12^+\) & \(\mu_n\)
& \(2.617\times10^{-10}\,\mathrm{s}\)\\

\(\Sigma^+\) & \(uus\) & \(1189.37\) & \(+1\) & \(-1\)
& \(\frac12^+\) & \(\mu_n-\mu_e\)
& \(8.018\times10^{-11}\,\mathrm{s}\)\\

\(\Sigma^0\) & \(uds\) & \(1192.64\) & \(0\) & \(-1\)
& \(\frac12^+\) & \(\mu_n\)
& \(7.4\times10^{-20}\,\mathrm{s}\)\\

\(\Sigma^-\) & \(dds\) & \(1197.45\) & \(-1\) & \(-1\)
& \(\frac12^+\) & \(\mu_n+\mu_e\)
& \(1.479\times10^{-10}\,\mathrm{s}\)\\
\hline
\(\Xi^0\) & \(uss\) & \(1314.86\) & \(0\) & \(-2\)
& \(\frac12^+\) & \(\mu_n\)
& \(2.90\times10^{-10}\,\mathrm{s}\)\\

\(\Xi^-\) & \(dss\) & \(1321.71\) & \(-1\) & \(-2\)
& \(\frac12^+\) & \(\mu_n+\mu_e\)
& \(1.639\times10^{-10}\,\mathrm{s}\)\\
\hline
\(\Omega^-\) & \(sss\) & \(1672.45\) & \(-1\) & \(-3\)
& \(\frac32^+\) & \(\mu_n+\mu_e\)
& \(8.21\times10^{-11}\,\mathrm{s}\)\\
\hline
\end{tabular}%
}
\caption{Proprietà degli iperoni fondamentali.
Tutti hanno numero barionico \(B=1\).
\(S\) indica la stranezza, \(Q\) la carica elettrica.}
\end{table}
\end{frame}

\begin{frame}{Iperoni: decadimenti principali}
    \begin{table}
\centering
\scriptsize
\begin{tabular}{ll}
\hline
\textbf{Iperone} & \textbf{Decadimenti principali nel vuoto (BR)}\\
\hline
\(\Lambda^0\)
& \(\Lambda^0\to p+\pi^-\) (\(64.1\%\))\\
& \(\Lambda^0\to n+\pi^0\) (\(35.9\%\))\\
\hline
\(\Sigma^+\)
& \(\Sigma^+\to p+\pi^0\) (\(51.5\%\))\\
& \(\Sigma^+\to n+\pi^+\) (\(48.4\%\))\\

\(\Sigma^0\)
& \(\Sigma^0\to\Lambda^0+\gamma\) (\(\sim100\%\))\\

\(\Sigma^-\)
& \(\Sigma^-\to n+\pi^-\) (\(99.85\%\))\\
\hline
\(\Xi^0\)
& \(\Xi^0\to\Lambda^0+\pi^0\) (\(99.52\%\))\\

\(\Xi^-\)
& \(\Xi^-\to\Lambda^0+\pi^-\) (\(99.89\%\))\\
\hline
\(\Omega^-\)
& \(\Omega^-\to\Lambda^0+K^-\) (\(67.7\%\))\\
& \(\Omega^-\to\Xi^0+\pi^-\) (\(24.3\%\))\\
& \(\Omega^-\to\Xi^-+\pi^0\) (\(8.6\%\))\\
\hline
\end{tabular}
\caption{Principali canali di decadimento degli iperoni.
I decadimenti sono deboli, eccetto quello
elettromagnetico della \(\Sigma^0\).}
\end{table}
\end{frame}

\subsection{Inner Core: Hyperon Matter}

\begin{frame}{Numeri quntici conservati in $\beta$-equilibrio: membri pesanti della famiglia.}
    \begin{columns}[T]
        \begin{column}{0.6\textwidth}
            \begin{itemize}
                \item Conservazione: Numero Barionico, Carica implica all'equilibrio $\mu_i=B_i\mu_B+Q_i\mu_Q$.
                \item $\beta$-equilibrio con fuga dei neutrini: equilibrio reazioni $\beta$ tipo $\Pneutron\to\Pproton\Pelectron\APnue$ impone $\mu_n=\mu_p+\mu_e$.
                \item Per specie non interagente $\mu_i=E_{F,i}^{tot}=\sqrt{k_{F,i}^2+m_i^2}$. La comparsa di una nuova particella diventa possibile per $\mu_i\geq E_i(0)$.
                \item Nuovi Leptoni oltre \Pelectron: \Pmuon, $m_{\Pmuon}=\SI{105.7}{\mega\ev}$, elettroni e muoni hanno stessi numeri quantici conservati $\mu_{\Pmuon}=\mu_e$, presenti per $\mu_e\geq m_{\Pmuon}$ e $n_p=n_e+n_{\Pmuon}$; \Ptauon, $m_{\Ptauon}=\SI{1776.9}{\mega\ev}$ non compare in modelli NS standard.
                \item Mesoni (\Pquark\APquark, stabili per interazioni forti, decadono attraverso interazioni deboli): \Pgpz, $m_{\Pgpz}=\SI{135.0}{\mega\ev}$, \Pgppm $m_{\Pgppm}=\SI{139.6}{\mega\ev}$, \PKpm, $m_{\PKpm}=\SI{493.7}{\mega\ev}$, \PKz, $m_{\PKz}=\SI{497.6}{\mega\ev}$; per i mesoni negativi $B_{\Pgpm}=B_{\PKm}=0$, $Q_{\Pgpm}=Q_{\PKm}=-1$ da cui $\mu_e=\mu_{\PKp}=\mu_{\Pgpm}$. Mesoni sono Bosoni: eccitazione modi mesonici $\omega_{\Pgpm}(k=0)=\mu_e$, $\omega_{\PKm}(k=0)=\mu_e$.
                \item Iperone $\Lambda$: $m_{\Lambda}=\SI{1116}{\mega\ev}$, \SI{176}{\mega\ev} pi\'u pesanti di \Pneutron, $B_{\Lambda}=1$, $Q_{\Lambda}=0$, equilibrio chimico: $\mu_{\Lambda}=\mu_n$. Soglia $\mu_n\approx m_{\Lambda}$ trascurando interazioni. Heavier hyperons: $S=-1$: $\Sigma$ $m_{\PSigma}\approx\SI{1190}{\mega\ev}$, $S=-2$: $m_{\PXizero,\PXiminus}=\num{1315},\num{1322}\si{\mega\ev}$, $S=-3$: \POmega.

            \end{itemize}
        \end{column}
        \begin{column}{0.4\textwidth}
            \begin{align*}
                &\Pneutron\leftrightarrow\Pproton+\Pelectron+\APnue\tag{$\mu_n=\mu_p+\mu_e$}\\
                &\Pneutron\leftrightarrow\Pproton+\Pgpm\tag{$\mu_n=\mu_p+\mu_{\Pgpm}$}\\
                &\PKm\leftrightarrow\Pmuon+\APnum\tag{$\mu_{\PKm}=\mu_{\Pmuon}$}\\
                &\Pneutron\leftrightarrow\Pproton+\Pmuon+\APnum\tag{$\mu_n=\mu_p+\mu_{\Pmuon}$}
            \end{align*}
            \begin{itemize}
                \item EOS Softening - La comparsa di nuove specie permette di redistribuire B e energia: a part\'a di densit\'a la pressione risulta inferiore rispetto a quella composta dalle specie piu leggere.
                \item Hyperon Puzzle: La comparsa di iperoni rende troppo soft per stabilit\'a EOS di NS con $M\geq2\msun{}$, che sono osservate; ma le interazioni repulsive possono cambiare la conclusione.
            \end{itemize}
        \end{column}
    \end{columns}

\end{frame}

\begin{frame}{Reazioni produzione muoni}
    \begin{table}
\centering
\scriptsize
\resizebox{\textwidth}{!}{%
\begin{tabular}{lll}
\hline
\textbf{Processo}
& \textbf{Reazione}
& \textbf{Soglia termodinamica}\\
\hline
Direct Urca
& \(n\to p+\mu^-+\bar\nu_\mu\)
& \(\mu_n-\mu_p=m_\mu c^2\)\\

Modified Urca
& \(n+N\to p+N+\mu^-+\bar\nu_\mu\)
& \(\mu_n-\mu_p=m_\mu c^2\)\\

Conversione leptonica
& \(e^-+N\to\mu^-+N+\nu_e+\bar\nu_\mu\)
& \(\mu_e=m_\mu c^2\)\\
\hline
\end{tabular}%
}
\caption{Processi deboli di produzione dei muoni nella materia
nucleare fredda, trasparente ai neutrini.
\(N=n,p\) indica un nucleone spettatore.
Le soglie indicano la comparsa termodinamica dei muoni,
non necessariamente l'apertura cinematica dei singoli processi.}
\end{table}
\end{frame}


\begin{frame}{Gas \(npe\) ideale in equilibrio beta}


\begin{columns}[T]
    \begin{column}{0.5\textwidth}
\textbf{Ipotesi:} \(T=0\), particelle non interagenti, neutrini non intrappolati, neutralità elettrica.
\[
\mu_n=\mu_p+\mu_e,\qquad
n_p=n_e \ \Rightarrow\ k_{F,p}=k_{F,e}.
\]

\textbf{Nucleoni non relativistici, elettroni ultrarelativistici}
\[
k_{F,N}\ll m_N,\qquad k_{F,e}\gg m_e.
\]
\vspace{-0.3cm}
\begin{align*}
\sqrt{m_n^2+k_{F,n}^2}
&=\sqrt{m_p^2+k_{F,e}^2}
 +\sqrt{m_e^2+k_{F,e}^2}\\
m_n+\frac{k_{F,n}^2}{2m_n}
&\simeq m_p+\frac{k_{F,e}^2}{2m_p}+k_{F,e}.
\end{align*}
    \end{column}
    \begin{column}{0.5\textwidth}
        Assumendo \(m_n\simeq m_p\), \(k_{F,e}\ll m_p\):
\[
\mu_e\simeq k_{F,e}
\simeq\frac{k_{F,n}^2}{2m_n}
=E_{F,n}^{(\mathrm{nr})}
\simeq56\,\mathrm{MeV}
\left(\frac{n_B}{n_0}\right)^{2/3}
\]
\[
\frac{n_e}{n_B}
\simeq\left(\frac{k_{F,e}}{k_{F,n}}\right)^3
\simeq\left(\frac{k_{F,n}}{2m_n}\right)^3
\simeq5.1\times10^{-3}\frac{n_B}{n_0}.
\]
    \end{column}
\end{columns}

\begin{table}
\centering
\scriptsize
\resizebox{\linewidth}{!}{%
\begin{tabular}{cccccc}
\hline
\(\boldsymbol{n_B/n_0}\)
& \(\boldsymbol{k_{F,n}}\)
& \(\boldsymbol{\mu_B}\)
& \(\boldsymbol{E_{F,n}^{(\mathrm{nr})}}\)
& \(\boldsymbol{\mu_e}\)
& \(\boldsymbol{n_e/n_B}\)\\
& [MeV] & [MeV] & [MeV] & [MeV] & \\
\hline
0.5 & 257 & 974  & 35  & 35  & \(2.6\times10^{-3}\)\\
1.0 & 324 & 994  & 56  & 54  & \(4.6\times10^{-3}\)\\
2.0 & 409 & 1024 & 89  & 82  & \(8.2\times10^{-3}\)\\
3.0 & 468 & 1050 & 117 & 105 & \(1.1\times10^{-2}\)\\
\hline
\end{tabular}%
}
\end{table}
\vspace{-0.3cm}
{\scriptsize
\(n_0=0.15\,\mathrm{fm}^{-3}\), \(\mu_B=\mu_n\).
La tabella riporta il calcolo del gas libero
senza le approssimazioni analitiche sopra utilizzate.
}
Il potenziale chimico degli elettroni pu\'o essere alto con densit\'a elettronica bassa. I muoni compaiono attorno a $n_B=3n_0$.
\end{frame}

\begin{frame}{Hyperon \PLambda, \PSigmaminus threshold}
    \begin{itemize}
        \item Critical condition for \PLambda onset, ignoring interactions: $\mu_{\Lambda}=m_{\Lambda}=\mu_n$ ($B=1$, $Q=0$). \PLambda pu\'o comparire quando diventa energeticamente favorevole sostituire neutrone che occupa stato vicino alla superficie di Fermi con \PLambda nel primo stato disponibile.
        \item In gas of free neutrons: $m_{\Lambda}=\mu_n=\sqrt{m_n^2+k_{F,c}^2}$ con $k_{F,c}=\sqrt{m_{\Lambda}^2-m_n^2}=\SI{602}{\mega\ev}$ critical Fermi momentum of neutrons, corresponding to critical density $n_c=\frac{1}{3\pi^2}k_{F,c}^3=\SI{0.96}{\per\cubic\femto\meter}=6.4n_0$.
        \item $m_{\PSigmaminus}=\SI{1197}{\mega\ev}>m_{\Lambda}$ ma $B_{\Sigma}=1$, $Q_{\Sigma}=-1$: $\mu_{\PSigmaminus}=\mu_B+\mu_e$ quindi la soglia \'e $m_{\Sigma}=\mu_B+\mu_e=\SI{1072}{\mega\ev}+\SI{125}{\mega\ev}$.
    \end{itemize}
    \begin{columns}[T]
        \begin{column}{0.55\textwidth}
\begin{table}
\centering
\small
\begin{tabular}{lcccc}
\hline
\textbf{Particella}
& \(\boldsymbol{n_c}\) [fm\(^{-3}\)]
& \(\boldsymbol{n_c/n_0}\)
& \(\boldsymbol{\mu_Q}\) [MeV]
& \(\boldsymbol{\mu_B}\) [MeV]\\
\hline
\(\mu^-\)    & 0.457 & 3.0 & -106 & 1050\\
\(\Sigma^-\) & 0.616 & 4.1 & -125 & 1072\\
\(\Lambda\)  & 1.26  & 8.4 & -123 & 1116\\
\hline
\end{tabular}
\caption{Densità critiche di comparsa di nuove particelle
nella materia stellare fredda, per un gas di particelle
non interagenti (\(n_0=0.15\,\mathrm{fm}^{-3}\)).
Nessun'altra specie compare fino a \(20n_0\).}
\end{table}
        \end{column}
        \begin{column}{0.45\textwidth}
            \begin{itemize}
                \item \PSigmaminus appare prima di \PLambda perch\'e contiene potenziale chimico elettrone.
                \item Il potenziale chimico dell'elettrone raggiunge il picco all'onset di \PSigmaminus a $\mu_e^{max}=\SI{125}{\mega\ev}$i: minore della massa del pione e dei mesoni pi\'u pesanti. La situazione cambia quando si tiene conto delle interazioni.
            \end{itemize}
        \end{column}
    \end{columns}
    \begin{itemize}
        \item Interactions have to be included to reproduce observed pulsar up to $2\msun$ but also appearence of $\PSigmaminus$ with nearly zero momentum from high Fermi momentum \Pelectron and \Pneutron, ie less momentum in the system result in smaller increase of pressure with baryon number density.
        \item 
    \end{itemize}
\end{frame}

\begin{frame}{Hyper-nuclei}
    \begin{columns}[T]
        \begin{column}{0.50\textwidth}
            \begin{itemize}
                \item Produzione: $\PKminus+^AZ\to^A_{\Lambda}Z\Pgpm$, $\Pgpp+^AZ\to^A_{\Lambda}Z+\PKplus$ - le reazioni conservano la stranezza e trasformano neutrone in iperone $\Lambda$.
                \item Potenziale di \PLambda in funzione della densit\'a: modello fenomenologico $U(n)=-a(\frac{n}{n_0})+b(\frac{n}{n_0})\expy{\gamma}$, primo termine attrattivo secondo repulsivo al crescere della densit\'a. Riscriviamo il potenziale imponendo che il potenziale abbia minimo in $n_0$: $U(n)=U_0[\frac{\gamma}{\gamma-1}\frac{n}{n_0}-\frac{1}{\gamma-1}(\frac{n}{n_0})\expy{\gamma}]$. Per $\gamma=2$: $U_{\PLambda}(n)=-30[2x-x^2]\si{\mega\ev}$ ($x=1$: $U_{\PLambda}=\SI{-30}{\mega\ev}$, $x=2$ $U_{\PLambda}=0$, $x=3$ $U_{\PLambda}(x)=-\SI{90}{\mega\ev}$).
                \item Quark counting rule: Nucleone $U_N\approx\SI{-50}{\mega\ev}$, $U_{\PLambda}\approx\SI{-30}{\mega\ev}$, $U_{\PSigma}\approx\SI{30+-20}{\mega\ev}$, $U_{\PXi}\approx\SI{-15}{\mega\ev}$, la comparsa di \PSigmaminus \'e favorita da $\mu_{\PSigmaminus}=\mu_n+\mu_{e}$ ma se interazione \'e repulsiva in materia nucleare la comparsa \'e ritardata/soppressa, per \PXi il potenziale attrattivo ne favorisce la comparsa rispetto al gas libero. Il potenziale attrattivo sembra diminuire col numero di quark non-strange (3 per nulceone, 2 per \PLambda e 1 per \PXi).
                \item Alcuni ipernuclei doppi come $^6_{\PLambda\PLambda}He$ suggeriscono debole attrazione \PLambda\PLambda.
            \end{itemize}
        \end{column}
        \begin{column}{0.50\textwidth}
            \begin{figure}[!ht]\includegraphics[trim={0cm 0cm 0 0},clip,keepaspectratio,width=0.90\textwidth]{Schaffner/BE_LambdaNuclei}\label{fig:Schaffner_BE_LambdaNuclei}\end{figure}
            \begin{itemize}
                \item $R=r_0A\expy{\frac{1}{3}}$ - energia cinetica di particella confinata in raggio R: $k=\frac{2\pi}{\lambda}=\frac{\pi}{R}$, $p=\hbar k$, $E_{kin}=\frac{p^2}{2m}\propto A\expy{-\frac{3}{2}}$ - livelli buca sferica: $V_0-B_{\PLambda}=\frac{\hbar^2 k^2}{2m_{\PLambda}}$, $B_{\PLambda}(A)\approx V_0-\frac{\pi^2\hbar^2}{2m_{\PLambda}r_0^2}A\expy{-\frac{2}{3}}$.
            \end{itemize}
        \end{column}
    \end{columns}
            \begin{itemize}
                \item curve: stati orbitali di \PLambda in potenziale nucleare (shell-model con potenziale Saxon-Woods) trascurando split spin-orbita.
                \item $A\to\infty$: energia di legame $B\approx\SI{30}{\mega\ev}$, quindi in materia nucleare uniforme $U_{\PLambda}(n_0)=-\SI{30}{\mega\ev}$.
            \end{itemize}
\end{frame}

\begin{frame}{Densit\'a critica per iperoni interagenti}
    \begin{itemize}
        \item Per \PLambda passiamo dalla condizione $\mu_{\PLambda}=\mu_n$ a $\mu_n=E_{\PLambda}(0;n)$ dove $E_{\PLambda}(0;n)=m_{\PLambda}+U_{\PLambda}(n)$ per primo \PLambda $k_{F,\PLambda}=0$ e il potenziale chimico neutronico \'e approssimato a $\mu_n(n)=\sqrt{k_{F,n}^2+m_n^2}+U_n(n)$, cio\'e $\sqrt{k_{F,n}^2+m_n^2}+U_n(n)=m_{\PLambda}+U_{\PLambda}(n)$.
        \item Conta la differenza tra i potenziali: $\Delta U_n^{crit}(n)=U_n(n)-U_{\PLambda}(n)=m_{\PLambda}-\sqrt{k_{F,c}^2+m_n^2}$: se aumentando densit\'a potenziale neutronico \'e pi\'u repulsivo di quello del \PLambda, $U_n(n)>U_{\PLambda}(n)$, pu\'o diventare favorevole sostituire parte dei neutroni con iperoni. Per $x=1$ i neutroni hanno $k_{F,n}\approx\SI{324}{\mega\ev}$ quindi per comparsa \PLambda $\Delta U_{crit}=\SI{123}{\mega\ev}$, per $x=2$ i neutroni hanno $k_{F,n}\approx\SI{410}{\mega\ev}$ quindi per comparsa \PLambda $\Delta U_{crit}=\SI{92}{\mega\ev}$, per $x=3$ i neutroni hanno $k_{F,n}\approx\SI{470}{\mega\ev}$ quindi per comparsa \PLambda $\Delta U_{crit}=\SI{66}{\mega\ev}$: pi\'u sale densit\'a meno deve contribuire la differenza dei potenziali per compensare maggiore massa di \PLambda.
        \item Per $\gamma=2$: $\Delta U(n)=U_n-U_{\PLambda}=(U_{n,0}-U_{\PLambda,0})(2x-x^2)=-20(2x-x^2)\si{\mega\ev}$ con $\Delta U(3n_0)=\SI{60}{\mega\ev}$ vicino alla formula fisica. Per $\gamma>2$ il contributo repulsivo neutronico cresce molto pi\'u rapidamente: incertezza su $\gamma$ produce incertezza sulla densit\'a di soglia.
        \item Energia di Simmetria: $\frac{E}{A}(n,\delta)=E_0(n)+S(n)\delta^2+\ldots$, $\delta=\frac{n_n-n_p}{n_B}$ quindi $\frac{E}{A}(n,1)\approx E_0(n)+S(n)$ con $S(n)>0$ pu\'o favorire la comparsa di nuove specie barioniche, e dato che $\mu_n=\TDy{n_n}{\epsilon}$ il contributo dell'energia di simmetria a $\mu_n$ dipende anche dalla sua derivata prima rispetto alla densit\'a.
        \item Quando compare \PSigma? Per gas libero $n_c^{\PSigmaminus}\approx4.1n_0$, $n_c^{\PLambda}\approx8.4n_0$; adesso sappiamo che $U_{\PLambda}(n_0)\approx\SI{-30}{\mega\ev}$, $U_{\PSigma}(n_0)\approx\SI{+30}{\mega\ev}$ e per le condizioni di comparsa: $\mu_n=U_{\PLambda}(n)+m_{\PLambda}$, $\mu_n+\mu_e=U_{\PSigmaminus}(n)+m_{\PSigmaminus}$. La repulsione nucleare aumenta energia di soglia per \PSigmaminus.
    \end{itemize}
\end{frame}

\begin{frame}{Hyperon Puzzle}
    \begin{itemize}
        \item \PLambda fells less repulsive potential then neutron-neutron: \PLambda becomes favored at high density $n\geq\numrange{2}{3}n_0$ (quark-counting rule: less hight repulsive potential).
        \item Nella materia nucleare simmetrica protoni e neutroni risentono di contributi isoscalari identici. La differenza energetica legata all'asimmetria protoni-neutroni \'e descritta dall'energia di simmetria $S(n_0)=\SI{32}{\mega\ev}$ molto piccola rispetto ai contributi scalari e vettoriali, centinaia di \si{\mega\ev}: in $\beta$-equilibrio la materia pu\'o rimanere neutron-rich con $\delta\approx1$.
        \item Sembra ragionevole che la repulsione isoscalare di \PLambda si inferiore a quella dei nucleoni: una volta raggiunta la soglia, aumentando $n_B$, diventa favorevole trasformare neutroni in \PLambda.
        \item Quark-counting rule in RMF: $g_{\omega\PLambda}=\frac{2}{3}g_{\omega N}$, il campo vettoriale $\omega$ contribuisce alla repulsione alle alte densit\'a.
        \item Se l'energia si distribuisce tra pi\'u specie la densit\'a di energia aumenta meno rapidamente con l'aumento della densit\'a: EOS diventa pi\'u soft.
        \item Modello con EOS dominato da interazioni $M_{max\propto yM_{Landau}}$, $y\gg1$ intensit\'a interazione repulsiva: $M_{max}^{\PLambda}\approx \frac{2}{3}M_{max}^{nucl}$. L'EOS diventa troppo soffice per sostenere una NS di $2\msun{}$ osservata.
        \item Soluzioni: ritardare comparsa di \PLambda $n_c^{\PLambda}>n_{max}^{star}$ ma contrasta con interazioni \PLambda\Pneutron attrattiva; repulsione YY pi\'u intensa ad alte densit\'a limita popolazione iperionica e rendere EOS pi\'u rigida e anche interazioni a 3 corpi NN\PLambda; prima che comparsi di \PLambda renda EOS troppo soffice la materia subisce transizione di fase con nuovi gradi di libert\'a, i quark.
    \end{itemize}
\end{frame}

\begin{frame}{$2\msun{}$ NS structure}
    \begin{columns}[T]
        \begin{column}{0.5\textwidth}
            \begin{itemize}
                \item Crosta Esterna(\SIrange{13.1}{12.9}{\kilo\meter},$\num{e-3}\epsilon_0$): lattice of nuclei in electron gas.
                \item Inner Crust (\SI{12.5}{\kilo\meter}, $0.5\epsilon_0$): lattice of nuclei in liquid of neutrons and electrons.
                \item Outer Core (\SI{9.8}{\kilo\meter}, $2\epsilon_0$): neutron matter with admixture of protons and electrons. $\epsilon_0$ at \SI{12}{\kilo\meter}.
                \item Inner Core ($3.6\epsilon_0$): hyperon may be present at $(2-3)\epsilon_0$, quark at $3\epsilon_0$ at \SI{5.9}{\kilo\meter}.
                \item GW from NS merger implies max radius \SI{13}{\kilo\meter}
                \item Upper bound on max energy density: causality. Stiffest EOS has $c_s=1$, and for $2\msun{}$ NS, $\epsilon_{max}\approx 13.5\epsilon_0$
            \end{itemize}
        \end{column}
        \begin{column}{0.5\textwidth}
            \begin{figure}[!ht]\includegraphics[trim={0cm 0cm 0 0},clip,keepaspectratio,width=0.95\textwidth]{Schaffner/NS_structure}\label{fig:Schaffner_NS_structure}\end{figure}
            \begin{itemize}
                \item Linee separano le diverse regioni, i colori le composizioni
            \end{itemize}
        \end{column}
    \end{columns}
    
\end{frame}

\section{Tre definizioni di massa nelle stelle di neutroni}
\subsection{Geometria e definizioni}

% 01
\begin{frame}{Tre nozioni di massa: il problema fisico}
\small
\begin{columns}[T,onlytextwidth]
\begin{column}{0.49\textwidth}
\textbf{Massa gravitazionale $M_G$}
\[
M_G=m(R).
\]
Determina il campo gravitazionale esterno;
e' quella misurata tramite le orbite.
\medskip
\textbf{Massa barionica $M_B$}
\[
M_B=m_bN_B.
\]
Misura il numero di barioni, moltiplicato
per una massa di riferimento $m_b$.
\end{column}
\begin{column}{0.49\textwidth}
\textbf{Massa propria $M_P$}
\[
M_P=\frac{1}{c^2}\int\varepsilon\,dV_{\rm prop}.
\]
Somma delle energie locali della materia,
integrate sul volume proprio.
\medskip
\textbf{Domanda centrale}
\par\medskip
Perche' $M_G$, $M_B$ e $M_P$ non coincidono
per una stella relativistica?
\end{column}
\end{columns}
\end{frame}

% 02
\begin{frame}{Geometria della stella: metrica TOV}
\small
\begin{columns}[T,onlytextwidth]
\begin{column}{0.49\textwidth}
Per una stella statica e sfericamente simmetrica:
\[
\begin{aligned}
ds^2={}&-e^{2\Phi(r)}c^2dt^2\\
&+\frac{dr^2}{1-2Gm(r)/(rc^2)}\\
&+r^2d\Omega^2.
\end{aligned}
\]
\[
d\Omega^2=d\theta^2+\sin^2\theta\,d\phi^2.
\]
\end{column}
\begin{column}{0.49\textwidth}
\textbf{Variabili locali}
\begin{itemize}
\item $\varepsilon(r)$: densita' di energia
nel riferimento comovente.
\item $n_B(r)$: densita' numerica barionica locale.
\item $m(r)$: funzione di massa gravitazionale.
\item $P(r)$: pressione isotropa.
\end{itemize}
\medskip
La superficie $R$ e' definita, nel modello ideale,
da $P(R)=0$.
\end{column}
\end{columns}
\end{frame}

% 03
\begin{frame}{Elemento di volume proprio: derivazione}
\small
\begin{columns}[T,onlytextwidth]
\begin{column}{0.49\textwidth}
Sulla superficie $t=\mathrm{cost.}$,
la distanza radiale propria e'
\[
dl_r=\frac{dr}{\sqrt{1-2Gm(r)/(rc^2)}}.
\]
L'area della sfera a raggio areale $r$ e'
\[
A(r)=4\pi r^2.
\]
Un guscio infinitesimo ha quindi volume
\[
dV_{\rm prop}=A(r)\,dl_r.
\]
\end{column}
\begin{column}{0.49\textwidth}
Sostituendo:
\[
\begin{aligned}
dV_{\rm prop}
&=4\pi r^2
\frac{dr}{\sqrt{1-2Gm(r)/(rc^2)}}.
\end{aligned}
\]
Per confronto, il volume coordinato e'
\[
dV_{\rm coord}=4\pi r^2dr.
\]
In generale:
\[
dV_{\rm prop}\geq dV_{\rm coord}.
\]
\end{column}
\end{columns}
\end{frame}

% 04
\begin{frame}{Massa gravitazionale: equazione TOV}
\small
Nel formalismo TOV, l'equazione per la funzione di massa e'
\[
\frac{dm}{dr}=\frac{4\pi r^2}{c^2}\,\varepsilon(r),
\qquad m(0)=0.
\]
Integrando dal centro alla superficie:
\[
\begin{aligned}
M_G\equiv m(R)
&=\int_0^R\frac{dm}{dr}\,dr\\
&=\frac{4\pi}{c^2}\int_0^R\varepsilon(r)r^2\,dr.
\end{aligned}
\]
\medskip
Qui compare $4\pi r^2dr$, \emph{non} l'elemento di volume proprio.
L'energia gravitazionale non e' descritta da una densita' locale
univoca da aggiungere a $\varepsilon$.
\end{frame}

% 05
\begin{frame}{Il significato osservativo di $M_G$}
\small
\begin{columns}[T,onlytextwidth]
\begin{column}{0.49\textwidth}
All'esterno della stella, $r\geq R$,
la soluzione e' Schwarzschild:
\[
\begin{aligned}
ds^2={}&-\left(1-\frac{2GM_G}{rc^2}\right)c^2dt^2\\
&+\frac{dr^2}{1-2GM_G/(rc^2)}\\
&+r^2d\Omega^2.
\end{aligned}
\]
L'unico parametro di massa esterno
per questo sistema non rotante e' $M_G$.
\end{column}
\begin{column}{0.49\textwidth}
\textbf{Interpretazione}
\begin{itemize}
\item Determina la geometria esterna.
\item In un sistema binario, controlla la dinamica orbitale.
\item Per una stella isolata statica asintoticamente piatta,
coincide con la massa ADM.
\end{itemize}
\medskip
Una stella indicata come $1.4M_\odot$
ha normalmente $M_G=1.4M_\odot$.
\end{column}
\end{columns}
\end{frame}

% 06
\begin{frame}{Numero barionico: la quantita' conservata}
\small
\begin{columns}[T,onlytextwidth]
\begin{column}{0.49\textwidth}
La corrente barionica e'
\[
J_B^\mu=n_Bu^\mu,
\]
dove $u^\mu$ e' la quadrivelocita'
del fluido.
\medskip
Se il numero barionico e' conservato:
\[
\nabla_\mu J_B^\mu=0.
\]
La carica corrispondente e' il numero totale
$N_B$ di barioni.
\end{column}
\begin{column}{0.49\textwidth}
Per la configurazione statica, su una
sezione spaziale ortogonale a $u^\mu$:
\[
N_B=\int n_B\,dV_{\rm prop}.
\]
Sostituendo il volume proprio:
\[
N_B=4\pi\int_0^R
\frac{n_B(r)r^2\,dr}
{\sqrt{1-2Gm(r)/(rc^2)}}.
\]
\medskip
$N_B$ non dipende dalla convenzione scelta
per esprimere una ``massa barionica''.
\end{column}
\end{columns}
\end{frame}

% 07
\begin{frame}{Massa barionica $M_B$: definizione}
\small
\[
M_B=m_bN_B
=4\pi m_b\int_0^R
\frac{n_B(r)r^2\,dr}
{\sqrt{1-2Gm(r)/(rc^2)}}.
\]
\begin{columns}[T,onlytextwidth]
\begin{column}{0.49\textwidth}
\textbf{Che cosa misura?}
\begin{itemize}
\item La quantita' di materia barionica,
nel senso del numero barionico totale.
\item La massa che si attribuirebbe a $N_B$
barioni liberi secondo la convenzione $m_b$.
\end{itemize}
\end{column}
\begin{column}{0.49\textwidth}
\textbf{Che cosa non misura?}
\begin{itemize}
\item Non e' l'integrale dell'energia totale locale.
\item Non comprende esplicitamente
l'energia cinetica, le interazioni e i leptoni.
\end{itemize}
\medskip
$m_b$ puo' essere la massa di un nucleone
o l'unita' di massa atomica.
\end{column}
\end{columns}
\end{frame}

% 08
\begin{frame}{Massa propria $M_P$: definizione}
\small
La massa propria integra tutta la densita' di energia
\emph{locale} della materia sul volume spaziale proprio:
\[
M_P=\frac{1}{c^2}\int_0^R\varepsilon(r)\,dV_{\rm prop}.
\]
Quindi
\[
M_P=\frac{4\pi}{c^2}\int_0^R
\frac{\varepsilon(r)r^2\,dr}
{\sqrt{1-2Gm(r)/(rc^2)}}.
\]
\begin{columns}[T,onlytextwidth]
\begin{column}{0.49\textwidth}
$\varepsilon$ include le masse a riposo,
l'energia cinetica e le interazioni
presenti nell'EOS.
\end{column}
\begin{column}{0.49\textwidth}
A differenza di $M_B$, $M_P$ non e'
semplicemente proporzionale al numero
di barioni.
\end{column}
\end{columns}
\end{frame}

\subsection{Confronto ed energie di legame}

% 09
\begin{frame}{Perche' $M_P\geq M_G$?}
\small
Mettiamo a confronto le definizioni:
\[
\begin{aligned}
M_G&=\frac{4\pi}{c^2}\int_0^R\varepsilon(r)r^2\,dr,\\[3pt]
M_P&=\frac{4\pi}{c^2}\int_0^R
\frac{\varepsilon(r)r^2\,dr}
{\sqrt{1-2Gm(r)/(rc^2)}}.
\end{aligned}
\]
Per una stella regolare, con $\varepsilon\geq0$ e senza orizzonte:
\[
0\leq\frac{2Gm(r)}{rc^2}<1
\quad\Longrightarrow\quad
\frac1{\sqrt{1-2Gm(r)/(rc^2)}}\geq1.
\]
Quindi
\[
M_P\geq M_G.
\]
\end{frame}

% 10
\begin{frame}{Limite di campo debole: energia gravitazionale}
\footnotesize
Quando $2Gm/(rc^2)\ll1$:
\[
\frac1{\sqrt{1-2Gm/(rc^2)}}
\simeq1+\frac{Gm}{rc^2}+\cdots.
\]
La differenza fra massa propria e gravitazionale diventa
\[
\begin{aligned}
M_P-M_G
&\simeq\frac{4\pi}{c^2}\int_0^R
\varepsilon(r)r^2\frac{Gm(r)}{rc^2}\,dr\\
&=\frac1{c^2}\int_0^R\frac{Gm(r)}{r}\,dm(r).
\end{aligned}
\]
In teoria newtoniana, l'energia potenziale di legame e'
\[
W=-\int_0^R\frac{Gm(r)}{r}\,dm(r).
\]
Pertanto $M_P-M_G\simeq |W|/c^2$ nel limite debole.
\end{frame}

% 11
\begin{frame}{Non esiste un ordinamento universale fra $M_B$ e $M_P$}
\small
\begin{columns}[T,onlytextwidth]
\begin{column}{0.49\textwidth}
Dalle definizioni:
\[
M_B-M_P=
\int_0^R\left[m_bn_B(r)-\frac{\varepsilon(r)}{c^2}\right]
dV_{\rm prop}.
\]
La differenza dipende dalla relazione
fra energia per barione e massa di riferimento:
\[
\frac{\varepsilon}{n_B}
\quad\hbox{contro}\quad m_bc^2.
\]
\end{column}
\begin{column}{0.49\textwidth}
L'energia locale contiene:
\begin{itemize}
\item energia di massa a riposo;
\item energia cinetica/degenerazione;
\item interazioni nucleari;
\item contributi dei leptoni.
\end{itemize}
\medskip
Non si deve assumere $M_B>M_P$
o $M_P>M_B$ senza conoscere l'EOS
e la scelta di $m_b$.
\end{column}
\end{columns}
\end{frame}

% 12
\begin{frame}{Energia di legame gravitazionale e totale}
\small
\begin{columns}[T,onlytextwidth]
\begin{column}{0.49\textwidth}
\textbf{Contributo gravitazionale}
\[
E_{\rm bind}^{\rm grav}
=(M_P-M_G)c^2.
\]
\medskip
Con questa definizione e' non negativo
per una stella regolare.
\medskip
E' associato alla differenza fra energia
locale integrata e massa gravitazionale
vista dall'esterno.
\end{column}
\begin{column}{0.49\textwidth}
\textbf{Energia di legame totale}
\[
E_{\rm bind}^{\rm tot}
=(M_B-M_G)c^2.
\]
\medskip
E' definita rispetto a $N_B$ barioni
isolati, ciascuno di massa $m_b$.
\medskip
Per le tipiche stelle di neutroni fredde
stabili, $M_B>M_G$.
\end{column}
\end{columns}
\end{frame}

% 13
\begin{frame}{Decomposizione dell'energia di legame}
\small
Aggiungiamo e sottraiamo $M_P$:
\[
\begin{aligned}
E_{\rm bind}^{\rm tot}
&=(M_B-M_G)c^2\\
&=(M_P-M_G)c^2+(M_B-M_P)c^2.
\end{aligned}
\]
\begin{columns}[T,onlytextwidth]
\begin{column}{0.49\textwidth}
\[
(M_P-M_G)c^2
\]
Contributo attribuito al legame
gravitazionale nella configurazione
statica considerata.
\end{column}
\begin{column}{0.49\textwidth}
\[
(M_B-M_P)c^2
\]
Differenza fra masse barioniche
di riferimento ed energia locale:
interazioni ed energia interna.
\end{column}
\end{columns}
\medskip
In relativita' generale la separazione dell'energia gravitazionale
in un contributo locale non e' univoca.
\end{frame}

\subsection{Applicazioni e riepilogo}

% 14
\begin{frame}{Due configurazioni con lo stesso numero di barioni}
\small
Confrontiamo due configurazioni stellari fredde, isolate e non rotanti:
\[
N_B^{(1)}=N_B^{(2)}
\quad\Longrightarrow\quad M_B^{(1)}=M_B^{(2)}.
\]
Le masse gravitazionali possono invece essere differenti:
\[
M_G^{(1)}\ne M_G^{(2)}.
\]
Per una trasformazione $1\to2$ con $M_G^{(1)}>M_G^{(2)}$,
la differenza disponibile come energia e'
\[
\Delta E=\bigl(M_G^{(1)}-M_G^{(2)}\bigr)c^2.
\]
\medskip
In un processo reale bisogna tenere conto dell'energia emessa,
delle eventuali perdite di barioni e delle altre quantita' conservate
(ad esempio il momento angolare).
\end{frame}

% 15
\begin{frame}{Calcolo di $M_G$, $M_B$, $M_P$: procedura TOV}
\small
Scelta una EOS $\varepsilon(P)$, si risolvono le equazioni TOV:
\[
\frac{dm}{dr}=\frac{4\pi r^2\varepsilon}{c^2},
\]
\[
\frac{dP}{dr}=
-\frac{G(\varepsilon+P)
\bigl[m+4\pi r^3P/c^2\bigr]}
{c^2r^2[1-2Gm/(rc^2)]}.
\]
Condizioni iniziali $m(0)=0$, $P(0)=P_c$;
la superficie e' definita da $P(R)=0$.
\medskip
Da $m(r)$, $\varepsilon(r)$ e $R$ otteniamo $M_G$ e $M_P$.
\medskip
Per calcolare $M_B$ serve anche la densita' barionica locale $n_B(P)$,
che non e' specificata direttamente dalla sola forma $P(\varepsilon)$.
\end{frame}

% 16
\begin{frame}{Ricostruire $n_B$ da un'EOS termodinamica}
\small
Per materia fredda, neutra e in equilibrio beta,
quando il numero barionico e' la sola carica conservata indipendente:
\[
d\varepsilon=\mu_B\,dn_B,
\qquad
\varepsilon+P=\mu_Bn_B.
\]
Sostituendo $\mu_B=(\varepsilon+P)/n_B$ nella prima:
\[
\begin{aligned}
d\varepsilon&=\frac{\varepsilon+P}{n_B}\,dn_B,\\
\frac{dn_B}{n_B}&=\frac{d\varepsilon}{\varepsilon+P(\varepsilon)}.
\end{aligned}
\]
Integrando:
\[
\ln\frac{n_B(\varepsilon)}{n_{B,*}}
=\int_{\varepsilon_*}^{\varepsilon}
\frac{d\varepsilon'}{\varepsilon'+P(\varepsilon')}.
\]
La costante $n_{B,*}$ si fissa tramite la normalizzazione
microfisica della EOS.
\end{frame}

% 17
\begin{frame}{Riepilogo: definizioni e limite newtoniano}
\footnotesize
\renewcommand{\arraystretch}{1.5}
\begin{tabular}{@{}p{0.17\textwidth}p{0.49\textwidth}p{0.25\textwidth}@{}}
\hline
\textbf{Massa} & \textbf{Definizione} & \textbf{Interpretazione}\\
\hline
$M_G$ & $\displaystyle\frac{1}{c^2}\int_0^R4\pi r^2\varepsilon\,dr$
& Campo esterno; massa osservata.\\
$M_B$ & $\displaystyle m_b\int_0^R n_B\,dV_{\rm prop}$
& Numero di barioni $\times m_b$.\\
$M_P$ & $\displaystyle\frac{1}{c^2}\int_0^R\varepsilon\,dV_{\rm prop}$
& Energia locale sul volume proprio.\\
\hline
\end{tabular}
\medskip
Per campo gravitazionale debole, se l'energia interna e le interazioni
sono piccole rispetto alla massa a riposo:
\[
M_G\simeq M_P\simeq M_B.
\]
Nelle stelle di neutroni relativistiche le differenze
sono fisicamente significative.
\end{frame}

